\documentclass[11pt]{article}
\usepackage[T1]{fontenc}
\usepackage{lmodern}
\usepackage[margin=1in]{geometry}
\usepackage{amsmath,amssymb,amsthm,mathtools}
\usepackage{enumitem,microtype,needspace,booktabs}
\usepackage[colorlinks=true,allcolors=blue,bookmarksnumbered=true]{hyperref}
\hypersetup{pdftitle={The equivariant cohomological Hikita conjecture for arbitrary quivers},pdfauthor={OpenAI}}
\newcommand{\C}{\mathbb C}
\newcommand{\Z}{\mathbb Z}
\newcommand{\Qbb}{\mathbb Q}
\newcommand{\Nbb}{\mathbb Z_{\geq0}}
\newcommand{\OO}{\mathcal O}
\newcommand{\KK}{\mathcal K}
\newcommand{\cR}{\mathcal R}
\DeclareMathOperator{\Gr}{Gr}
\DeclareMathOperator{\GL}{GL}
\DeclareMathOperator{\Hom}{Hom}
\DeclareMathOperator{\End}{End}
\DeclareMathOperator{\Lie}{Lie}
\DeclareMathOperator{\Spec}{Spec}
\DeclareMathOperator{\Sym}{Sym}
\DeclareMathOperator{\Hilb}{Hilb}
\DeclareMathOperator{\rk}{rk}
\DeclareMathOperator{\codim}{codim}
\DeclareMathOperator{\im}{im}
\DeclareMathOperator{\Res}{Res}
\DeclareMathOperator{\eu}{eu}
\DeclareMathOperator{\wt}{wt}
\newcommand{\id}{\mathrm{id}}
\theoremstyle{plain}
\newtheorem{theorem}{Theorem}[section]
\newtheorem{lemma}[theorem]{Lemma}
\newtheorem{proposition}[theorem]{Proposition}
\newtheorem{corollary}[theorem]{Corollary}
\theoremstyle{definition}
\newtheorem{definition}[theorem]{Definition}
\newtheorem{remark}[theorem]{Remark}
\numberwithin{equation}{section}
\setlist[enumerate]{itemsep=3pt,topsep=5pt}
\setlist[itemize]{itemsep=3pt,topsep=5pt}
\setlength{\emergencystretch}{2em}
\title{The equivariant cohomological Hikita conjecture\\for arbitrary quivers}
\author{OpenAI}
\date{September 24, 2026}
\begin{document}
\maketitle
\begin{abstract}
We prove the equivariant cohomological Hikita conjecture for arbitrary finite
quivers, including loops and multiple arrows. For any dimension and framing
vectors, any commuting flavor torus, and a stability character whose
semistable locus is stable and has free gauge action, the equivariant
cohomology of the Nakajima variety is canonically isomorphic to the coordinate
ring of the scheme-theoretic fixed locus of the stability cocharacter on the
flavor-deformed Coulomb branch. This is an isomorphism of graded algebras over
the common coefficient ring, retaining nilpotents and including the empty case.
\end{abstract}
\section{Introduction}
\label{sec:introduction}

Hikita's conjecture relates two different kinds of algebra: the cohomology
of a symplectic resolution and the functions on a fixed subscheme of a
corresponding Coulomb branch.  The scheme structure is essential.  Even
when the underlying fixed-point set is small, its nilpotent functions can
encode positive-degree cohomology.  For quiver gauge theories, the
tautological classes on the Higgs side and the equivariant parameters on
the Coulomb side give a common coefficient ring.  The question is whether
the relations among those coefficients agree.

We prove this agreement for arbitrary finite quivers, under the precise
regularity hypothesis stated below.  The proof keeps the coefficient maps
and the full fixed schemes throughout.  Its two main comparisons have
different purposes: a calculation for the full topological torus
identifies the cohomology algebra, while a separate argument identifies
that torus-fixed scheme with the fixed scheme of the stability
cocharacter.

\subsection{The statement and its coefficient maps}

Let $Q=(Q_0,Q_1)$ be a finite oriented quiver.  Loops and multiple arrows are
allowed.  Fix complex vector spaces $V_i,W_i$ of dimensions $v_i,w_i$, and put
\begin{equation}\label{eq:gauge-data}
 G=\prod_{i\in Q_0}\GL(V_i),\qquad
 N=\bigoplus_{a:i\to j}\Hom(V_i,V_j)
       \oplus\bigoplus_{i\in Q_0}\Hom(V_i,W_i).
\end{equation}
An algebraic torus $F$ acts linearly on $N$ and commutes with $G$; write
$\mathfrak f=\Lie F$.  The induced action on $T^*N=N\oplus N^*$ is Hamiltonian.
For $(s,p)\in N\oplus N^*$, our cotangent moment map is characterized by
$\langle\mu(s,p),\xi\rangle=\langle p,\xi s\rangle$.
Choose a character $\nu$ of $G$ for which the $\nu$-semistable locus in
$\mu^{-1}(0)$ equals the stable locus and $G$ acts freely there.  We call such
a character \emph{regular}, and set
\begin{equation}\label{eq:higgs-definition}
 D=\mu^{-1}(0)^{\nu\text{-stable}},\qquad X=D/G.
\end{equation}
Thus $X$ is a smooth Nakajima variety.  We allow $D$ and $X$ to be empty.

Put $\OO=\C[[t]]$ and $\KK=\C((t))$.  The Braverman--Finkelberg--Nakajima
space of triples and its flavor-deformed Coulomb algebra are
\begin{align}
 \cR&=\{[g,s]\in G(\KK)\times_{G(\OO)}N(\OO):gs\in N(\OO)\},
       \label{eq:triples-definition}\\
 A_F&=H_*^{G(\OO)\times F}(\cR;\C),\qquad
 Y_{\mathfrak f}=\Spec A_F.\label{eq:coulomb-definition}
\end{align}
Here $(g,s)\sim(gh^{-1},hs)$, and the multiplication is BFN convolution.
We use the commutative algebra, without the loop-rotation deformation
parameter.  Its topological grading is by
$\pi_1(G)=\Z^{\{i:v_i>0\}}$, and defines a torus
$T_{\mathrm{top}}=\Hom(\pi_1(G),\C^\times)$.
The character $\nu$ gives a cocharacter of this torus.  Its fixed locus is
scheme-theoretic: if $(A_F)_m$ denotes its weight-$m$ subspace, then
\begin{equation}\label{eq:nu-quotient}
 \C[Y_{\mathfrak f}^{\nu}]
       =A_F/\langle(A_F)_m:m\ne0\rangle.
\end{equation}
No reduction of this quotient will be taken.

Let $T\subset G$ be the diagonal maximal torus, let $\mathcal W$ be its Weyl
group, and write
\begin{equation}\label{eq:coefficient-ring}
 R=H^*_{G\times F}(\mathrm{pt};\C)
   =\C[\Lie T\oplus\mathfrak f]^{\mathcal W},\qquad
 R_0=H_G^*(\mathrm{pt};\C).
\end{equation}
The Kirwan map $\phi_H:R\to H_F^*(X;\C)$ and the coefficient map
$\phi_C:R\to\C[Y_{\mathfrak f}^{\nu}]$ are the natural maps on the two
sides of the Hikita conjecture.  On the first side, a gauge Chern class
maps to the Chern class of its tautological bundle, and a flavor parameter
acts through $H_F^*(\mathrm{pt})$.  On the second side, $R$ acts by the
coefficient classes of the Coulomb algebra, followed by the fixed-scheme
quotient.  Each linear gauge or flavor parameter has degree two.

\begin{theorem}\label{thm:main}
For every datum \eqref{eq:gauge-data}, every commuting flavor torus $F$, and
every regular stability character $\nu$, both coefficient maps are
surjective and have the same kernel.  They therefore induce a canonical
isomorphism of graded
$R$-algebras
\[
 H_F^*(X;\C)\ \cong\ \C[Y_{\mathfrak f}^{\nu}].
\]
The gradings are ordinary cohomological degree and BFN homological degree,
respectively.  The assertion includes the empty case, whose ring is zero.
\end{theorem}

This is the equivariant cohomological Hikita conjecture for quiver gauge
theories, in the form of \cite[Conjecture~2.1, Equation~(2.3)]{DK25}.
It concerns the full coefficient-algebra quotient, with its nilpotents.
No finite-type, Dynkin, conical, or good/ugly hypothesis is imposed.

\subsection{History and the comparison problem}

A classical model is the description of type-$A$ Springer-fiber
cohomology through scheme-theoretic intersections of nilpotent orbit
closures with a Cartan subalgebra, developed by de Concini--Procesi and
Tanisaki \cite{DeConciniProcesi81,Tanisaki82}.  In the setting of
symplectic duality \cite{BLPW16}, Hikita formulated a fixed-scheme
realization of cohomology and established it for the Hilbert scheme of
points in the plane, with related type-$A$ Spaltenstein and smooth
hypertoric examples \cite{Hikita17}.  These examples explain why the
infinitesimal structure of the fixed locus, rather than its set of
points, is central to the question.

Nakajima proposed a mathematical Coulomb-branch construction for
three-dimensional gauge theories \cite{Nakajima16};
Braverman--Finkelberg--Nakajima constructed its commutative convolution
algebra \cite{BFN18,BFN19}, extending the affine-Grassmannian convolution
framework of Bezrukavnikov--Finkelberg--Mirkovi\'c \cite{BFM05}.
For quiver gauge theories the corresponding
Higgs quotient is a Nakajima variety \cite{Nakajima94}.
The coefficient-compatible formulation has important predecessors.
Kamnitzer--Tingley--Webster--Weekes--Yacobi record Nakajima's equivariant
conjecture in the affine-Grassmannian setting and strengthen its
specialized form by requiring compatibility with tautological classes
\cite[Conjectures~8.9--8.10]{KTWWY19}.
Hilburn--Kamnitzer--Weekes formulate the equivariant quantized refinement
over the common coefficient algebra
\cite[Section~5.7, Conjecture~5.26]{HKW23}.
Krylov--Shlykov state the general-quiver comparison for the
\emph{full topological-torus fixed scheme} as equality of coefficient-map
kernels \cite[Section~8.1, Conjecture~8.1]{KrylovShlykov25}.
We first establish that equality and then compare the full-torus fixed
scheme with the stability-cocharacter fixed scheme.
Thus compatibility with the coefficient maps is part of the conjectural
formulation itself.

Krylov--Shlykov prove the Gieseker, or framed Jordan-quiver, case
\cite{KrylovShlykov25}.  Dumanski--Krylov give the formulation used here
and compare cohomological, equivariant, and $K$-theoretic versions
\cite{DK25}.  Their finite-type conclusion includes types $A,D,E_6$;
in types $E_7,E_8$ it imposes the framing restrictions in
\cite[Corollary~2.17]{DK25}.  Their Remark~2.18 explains why the earlier
unrestricted ADE assertion must be qualified.
Recent developments also include Setiabrata's Hikita surjectivity result
for the base-affine-cotangent example \cite{Setiabrata24}, and the refined quantum
Hikita theorem of Dinkins--Karpov--Krylov for ADE quivers with
minuscule-supported framing and for the Jordan quiver
\cite[Theorem~B]{DinkinsKarpovKrylov26}.
Theorem~\ref{thm:main} treats the cohomological algebra for all the finite
quiver data above.  It does not assert the quantized statement in
\cite[Equation~(2.2)]{DK25}, nor extend the theorem to arbitrary
reductive gauge groups; outside this setting, proposed Hikita-type
identifications can fail \cite[Example~6.7]{HoangKrylovMatvieievskyi24}.

Several established tools enter the argument in distinct ways.
BFN abelianization turns minuscule convolution classes into explicit
Laurent expressions, and Weekes's generation theorem makes these classes
sufficient for arbitrary quivers \cite{BFN18,BFN19,Weekes19}.
McGerty--Nevins's ordinary and equivariant Kirwan surjectivity results
supply the tautological presentation on the Higgs side
\cite[Theorem~1.2 and Corollary~1.5]{MN18}.
Hausel's arithmetic Fourier-transform formula computes the Betti numbers
that eventually determine the size of the coefficient quotient
\cite{Hausel06}.  Character-linearized GIT and the topological
Kempf--Ness correspondence connect regular stability to the flat
hyperk\"ahler quotient construction
\cite{King94,Hoskins14,HKLR87,Nakajima94}.
Our use of these inputs is explained at their precise interfaces below.

The finite-dimensional critical-model viewpoint is related to Nakajima's
original construction and to Botta--Tamagni's critical lattice and
quasimap models \cite{Nakajima16,BottaTamagni26}.  Here the specific
stable-pair test proves an inclusion of coefficient ideals; it does not
require their entire Coulomb-action construction.  Deformation arguments
also appear in Krylov--Shlykov's Gieseker proof
\cite[Remark~1.15]{KrylovShlykov25}.  Our loop-removal step uses subgroup
localization \cite{Brion98}, with the prequotient perspective of
\cite[Proposition~2.3]{DK25}, and an explicit descent for homogeneous
ideals.  These connections distinguish the established methods from the
particular comparison arguments carried out here.

The difficulty is not just to match dimensions.  One must identify the
relations under the specified maps, with their grading and nilpotents,
for arbitrary flavors and for every regular stability.  Neither a
calculation on closed fixed points nor an isomorphism after inverting
generic parameters establishes that assertion.  A further distinction is
that a regular character can satisfy $\nu\cdot\mathbf b=0$ for an integer
vector $0<\mathbf b\leq\mathbf v$.  Therefore the full topological-torus calculation
does not by itself identify the cocharacter-fixed quotient.

\subsection{Proof strategy and organization}

Define the intermediate coefficient algebra and the two ideals to be
compared by
\begin{equation}\label{eq:intro-two-ideals}
 \begin{split}
 B_F&=\C[Y_{\mathfrak f}^{T_{\mathrm{top}}}]
      =A_F/\langle A_F^{(d)}:d\ne0\rangle,\\
 I_{\mathrm{top}}&=\ker(R\longrightarrow B_F),\qquad
 I_H=\ker(R\xrightarrow{\phi_H}H_F^*(X)),
 \end{split}
\end{equation}
where $A_F^{(d)}$ is the subspace of topological charge
$d\in\pi_1(G)$; its cocharacter weight is $\nu\cdot d$.  Neither of
these is the cohomological degree used in Hilbert series.
Section~\ref{sec:preliminaries} explains why the
coefficient maps to $B_F$ and to $H_F^*(X)$ are surjective.
Thus the first task is precisely the equality
$I_{\mathrm{top}}=I_H$.  The second task is the equality of the
cocharacter and full-torus defining ideals inside $A_F$.
A subscript $0$ means that the flavor torus is trivial, so
$R_0=H_G^*(\mathrm{pt})$ and $B_0$ denotes the corresponding full-torus
quotient.  Figure~\ref{fig:proof-route} separates the two tasks.

\paragraph{1. Stable pairs give a quotient map.}
Sections~\ref{sec:correspondences} and \ref{sec:stable-pairs} prove
$I_{\mathrm{top}}\subseteq I_H$ for all the data of the theorem.
The nontrivial minuscule generators used here are dressed lattice
modifications whose coweight entries are all in $\{0,1\}$ or all in
$\{0,-1\}$, with nonzero topological charge; a subset may involve several
vertices.  A convolution word is a finite product of these classes.
By minuscule generation, a coefficient in $I_{\mathrm{top}}$ is a finite
sum of positive-length convolution words, each of total topological charge zero.
Section~\ref{sec:correspondences} realizes each such word by
finite-dimensional lattice correspondences, keeping all degree shifts.
The test in Section~\ref{sec:stable-pairs} evaluates them on stable
cotangent pairs.  The stability character assigns an additive degree to
each modification, so the degrees in a neutral word sum to zero.
After reordering each word, its first modification is nontrivial and has
nonpositive character-line degree; stability excludes it.  On the identity
correspondence, the same test gives a split restriction onto $H_F^*(X)$,
which detects the coefficient class itself.  The tested relation therefore vanishes
in cohomology, yielding the graded coefficient-compatible surjection
$B_F\twoheadrightarrow H_F^*(X)$ of
Proposition~\ref{prop:kernel-inclusion}.

\paragraph{2. Added loops make the size comparison explicit.}
Temporarily suppose that the quiver has a loop at every vertex of
positive gauge dimension, take no flavor, and use determinant stability
$\nu_{\det}=\prod_i\det_{V_i}$.
The loop weights cancel the Grassmannian denominators in minuscule
abelianization.  Section~\ref{sec:laurent} proves a Laurent model with
polynomial coefficients and filters $B_0$ by multipartitions
of the gauge dimension vector.  Its successive quotients are bounded by
invariant parts of polynomial quotients defined by cut products: products of framing variables
and differences across a subset of slots and its complement.
Section~\ref{sec:counting} computes their Hilbert series
using a reachability stratification and identifies the resulting sum
with Hausel's ordinary Poincar\'e polynomial.  The upper bound and the
surjection from the first step force equality in every degree.  This is
an isomorphism under the original coefficient map, not merely a numerical
identity; see Corollary~\ref{cor:loop-isomorphism}.

\paragraph{3. Deformation returns to the original data.}
For an arbitrary quiver $Q$, let $Q^+$ be obtained by adding one loop at
each vertex of positive gauge dimension.  Give the new loops the common
weight one of an auxiliary flavor torus $H=\C^\times$, with parameter
$z\in H_H^2(\mathrm{pt})$.
The comparison in Section~\ref{sec:deformations} uses a lift of the
no-flavor isomorphism for $Q^+$ to this flavor.  Near zero gauge parameters and nonzero mass
$z=1$, the new minuscule Euler factors are invertible on the Coulomb
side.  Subgroup localization on the stable prequotient gives the
corresponding comparison on the Higgs side.  The resulting equality of
localized coefficient ideals descends to the original homogeneous ideals
in $R_0$: their equality at the graded origin determines them globally.
This homogeneity argument, not specialization of an unspecified generic
isomorphism, is the content of Proposition~\ref{prop:loop-removal}.
Local constancy of regular stability, quaternionic rotation, and generic
constancy of compactly supported Betti numbers then compare all regular
characters.  A final flavor lift proves
$B_F\simeq H_F^*(X)$ in Theorem~\ref{thm:full-torus}.

\paragraph{4. Complementary monopoles identify the cocharacter quotient.}
Section~\ref{sec:cocharacters} proves, independently of the preceding
cohomology comparison, that the natural map
$\C[Y_{\mathfrak f}^{\nu}]\to B_F$ is an isomorphism.
It combines polynomial-form Hall-algebra generation
\cite{JindalNegut26} with the monopole maps of
Lemma~\ref{lem:polynomial-monopole-maps}, derived from
\cite[Theorem~5.7]{Tsymbaliuk23Difference} with explicit product signs.
The resulting finite identities can be specialized without dividing by
rotation or arrow parameters.
After adjoining a rank-one framing vertex, any surviving nontrivial minuscule
monopole has a surviving complementary monopole.  Hall generation would
then decompose both charges into indecomposable charges of the enlarged
quiver, all of weight zero for the extended character.
Lemma~\ref{lem:root-moment} realizes these charges by moment-level blocks.
Their semistable direct sum has a stabilizer that remains
positive-dimensional even after removing the ineffective diagonal
scalar.  This contradicts the freeness in regularity, including when
the semistable locus is empty.  Proposition~\ref{prop:cocharacter-fixed}
therefore identifies the two defining ideals, completing the
scheme-theoretic comparison.

\begin{figure}[hbp]
\centering
\begingroup
\small
\setlength{\fboxsep}{7pt}
\begin{minipage}[t]{0.47\textwidth}
\centering
\textbf{The coefficient-ideal comparison}\par\medskip
\fbox{\parbox{0.89\linewidth}{\centering
Every $Q,F$ and regular $\nu$\\
Stable-pair test\\
$I_{\mathrm{top}}\subseteq I_H$\\
Sections~\ref{sec:correspondences}--\ref{sec:stable-pairs}}}
\par\smallskip$\Downarrow$\par\smallskip
\fbox{\parbox{0.89\linewidth}{\centering
$Q^+$, no flavor, determinant stability\\
Laurent model and cut-ideal count\\
$B_0(Q^+)\simeq H^*(X_{\det}(Q^+))$\\
Sections~\ref{sec:laurent}--\ref{sec:counting}}}
\par\smallskip$\Downarrow$\par\smallskip
\fbox{\parbox{0.89\linewidth}{\centering
Lift to the added-loop flavor $H$\\
Compare near gauge $0$, mass $z=1$\\
Descend homogeneous ideals for $Q$\\
Proposition~\ref{prop:loop-removal}}}
\par\smallskip$\Downarrow$\par\smallskip
\fbox{\parbox{0.89\linewidth}{\centering
All regular stabilities, then all flavors\\
$B_F\simeq H_F^*(X)$\\
Theorem~\ref{thm:full-torus}}}
\end{minipage}
\hfill
\begin{minipage}[t]{0.47\textwidth}
\centering
\textbf{The fixed-scheme comparison}\par\medskip
\fbox{\parbox{0.89\linewidth}{\centering
Suppose a nontrivial minuscule class\\
survives in the cocharacter quotient.\\
Adjoin a rank-one framing vertex;\\
its complementary class also survives}}
\par\smallskip$\Downarrow$\par\smallskip
\fbox{\parbox{0.89\linewidth}{\centering
Polynomial Hall generation\\
Both surviving charges decompose\\
into indecomposable charges of the\\
enlarged quiver, all of weight zero\\
Lemma~\ref{lem:root-charge-generation}}}
\par\smallskip$\Downarrow$\par\smallskip
\fbox{\parbox{0.89\linewidth}{\centering
Lemma~\ref{lem:root-moment} gives moment-level blocks.\\
Their direct sum retains a\\
positive-dimensional stabilizer after\\
removing the diagonal scalar,\\
contradicting regularity}}
\par\smallskip$\Downarrow$\par\smallskip
\fbox{\parbox{0.89\linewidth}{\centering
$\C[Y_{\mathfrak f}^{\nu}]\simeq B_F$\\
as quotients of $A_F$\\
Proposition~\ref{prop:cocharacter-fixed}}}
\end{minipage}
\par\bigskip
$\displaystyle
 \C[Y_{\mathfrak f}^{\nu}]\ \xrightarrow{\ \sim\ }\
 B_F\ \xrightarrow{\ \sim\ }\ H_F^*(X)$
\endgroup
\caption{The two comparisons in the proof.  Here $Q^+$ adds one loop at
each positive-rank vertex of $Q$; dimension and framing vectors are
unchanged and suppressed.  The left column uses the stable-pair
surjection in the dimension comparisons.  The right column is
independent of the left and identifies defining ideals, not just closed
fixed points.  Both final maps preserve the coefficient ring and grading.}
\label{fig:proof-route}
\end{figure}


\section{Coefficient rings, minuscule classes, and flavor}
\label{sec:preliminaries}

This section supplies the generation and coefficient maps needed for the
two ideal comparisons described in Section~\ref{sec:introduction}.  Minuscule
generation makes $B_F$ a quotient of $R$; Kirwan surjectivity does the same
for $H_F^*(X)$.  Flavor base change and local constancy of regular
stability will let us apply these presentations to the auxiliary data
used later.

We use complex coefficients throughout.  Vertices with $v_i=0$ can be
discarded.  If all $v_i$ vanish, both sides of Theorem~\ref{thm:main} equal
$H_F^*(\mathrm{pt})$, so we henceforth allow ourselves to omit that case.
Write
\[
 P_F=\C[x_{ia}:i\in Q_0,\ 1\le a\le v_i]\otimes\C[\mathfrak f],
 \qquad \mathcal W=\prod_i S_{v_i},\qquad R=P_F^{\mathcal W}.
\]
The $x_{ia}$ are the Chern roots of the gauge bundles.  Each linear
form in $P_F$ has cohomological degree two.  An element
$\lambda=(\lambda_{ia})\in X_*(T)$ has topological charge
$[\lambda]=(\sum_a\lambda_{ia})_i$.
We write $A_F^{(d)}$ for the subspace of charge $d$ and define
\begin{equation}\label{eq:full-fixed-quotient}
 B_F=A_F/\langle A_F^{(d)}:d\ne0\rangle.
\end{equation}
Thus $B_F$ is the coordinate ring of the fixed subscheme for
$T_{\mathrm{top}}$.  A subscript $0$ means that the flavor torus is trivial.

\begin{lemma}\label{lem:trivial-summands}
For the purposes of Theorem~\ref{thm:main}, one may assume that the flavor
action assigns a character to each arrow summand and acts diagonally on
each framing multiplicity space.  Adding or removing a $G$-trivial
representation with any flavor action preserves both coefficient-algebra
quotients and all stability conditions.
\end{lemma}

\begin{proof}
Decompose $N$ into $G$-isotypic summands.  A commuting torus acts on each
multiplicity space and is diagonalizable there.  The nontrivial irreducible
gauge representations that occur are the bifundamentals between distinct
vertices, the dual standard representations from framing, and the
traceless endomorphisms at vertices of rank at least two.  All remaining
summands are gauge-trivial.  A weighted traceless endomorphism summand can
be replaced by the full endomorphism representation by adding a scalar
line of the same flavor weight.  This gives the asserted description.

If $E$ is gauge-trivial, the moment map for $N\oplus E$ is the moment map
for $N$, and semi-invariant generation shows that the stable and
semistable loci are the corresponding loci for $N$, times $T^*E$.
Their stabilizers are unchanged.  The quotient gains a factor $T^*E$,
which retracts $F$-equivariantly to zero.  In the space of triples the
condition on the $E(\OO)$ coordinate imposes no lattice constraint, so
this coordinate gives a common vector-space factor in every finite
truncation.  The normalized equivariant Thom isomorphism removes this
factor, compatibly with convolution, both gradings, and the coefficient
map.  This proves the claim on both sides.
\end{proof}

We call a coweight \emph{positive minuscule} if all its entries, across
all vertices, are $0$ or $1$, and \emph{negative minuscule} if all entries
are $0$ or $-1$.  Its orbit $\Gr_G^\lambda$ in
$\Gr_G=G(\KK)/G(\OO)$ is a product of ordinary Grassmannians and is
closed and smooth.

Write $\mathcal W_\lambda=\operatorname{Stab}_{\mathcal W}(\lambda)$,
$u_+=\max(u,0)$, and let $\wt(N)$ denote the multiset of weights of $N$.
For $p\in P_F^{\mathcal W_\lambda}$, put $p_{w\lambda}=w(p)$; stabilizer
invariance makes this independent of the representative $w$.

\begin{proposition}[Abelianization and minuscule generation]
\label{prop:minuscule}
There is an injective coefficient-preserving algebra map from $A_F$ into
\[
 \bigl(\operatorname{Frac}(P_F)[X_*(T)]\bigr)^{\mathcal W},\qquad
 e^\lambda e^\gamma=e^{\lambda+\gamma}.
\]
Besides $R$, the algebra is generated by the dressed classes associated
with all nonzero positive and negative minuscule coweights.  For
$p\in P_F^{\mathcal W_\lambda}$, the class has expression
\begin{equation}\label{eq:minuscule}
 \sum_{\lambda'\in\mathcal W\lambda}
 e^{\lambda'}p_{\lambda'}
 \frac{\prod_{\beta\in\wt(N)}
           \beta^{(-\langle\beta,\lambda'\rangle)_+}}
      {\eu_T(T_{t^{\lambda'}}\Gr_G^\lambda)}.
\end{equation}
Each distinct translate occurs once, and the dressing is translated
with the coweight.  Denote this class by $M_{\lambda,p}$.
Weights $\beta$ include their flavor components, but
pairings use only the gauge components; factors with exponent zero are
omitted.  The class has charge $[\lambda]$.  On charge zero, the BFN
grading is twice the ordinary polynomial degree of the rational
coefficients in this description.
\end{proposition}

\begin{proof}
We use the classical BFN localization and flavor construction
\cite[Proposition~3.22, Lemmas~5.9--5.13, Proposition~5.15,
and Corollary~5.21]{BFN18}, and its minuscule formula
\cite[Proposition~6.6]{BFN18}\cite[Appendix~A, Proposition~A.2]{BFN19},
together with quiver minuscule generation
\cite[Proposition~3.1 and Remarks~3.3--3.4]{Weekes19}.
The form of the formula can be checked directly: the inclusion of the
space of triples into the bundle of lattices $gN(\OO)$ contributes the
Euler class of
\[
 t^{\lambda'}N(\OO)\big/
       \bigl(N(\OO)\cap t^{\lambda'}N(\OO)\bigr),
\]
which is the numerator of \eqref{eq:minuscule}.  Localization on the
closed minuscule orbit gives the tangent denominator.  There is no
rotation, so translating a modification by $t^\eta$ changes the Laurent
monomial but does not shift the coefficients.  The localization theorem
and the orbit filtration give injectivity before taking Weyl invariants.
Equivalently, the image of the nonabelian algebra lies in the invariant
part of the displayed ring.

For completeness, the minuscule coweights used in the generation theorem
are subsets of slots at several vertices at once.  Decomposing positive
or negative integral coweights by successive level sets gives exactly
these two families; no assertion that single-slot monopoles suffice is
needed.  Lemma~\ref{lem:trivial-summands} puts arbitrary commuting flavors
in weighted quiver form.  Flavor characters change the equivariant Euler
factors but not the gauge-weight hyperplane arrangement used in the
generation proof; the same chamber argument therefore retains those
parameters, as discussed in \cite[Section~1.6.1]{Weekes19}.
The embedding used here is the classical localization map, not a
specialization of a representation in which loop rotation is inverted.
Finally, put $c_\lambda=\sum_{\beta\in\wt(N)}
(-\langle\beta,\lambda\rangle)_+$ and
$a_\lambda=\dim\Gr_G^\lambda$.  For homogeneous $p$, the BFN homological
normalization gives
\[
 \deg M_{\lambda,p}=2c_\lambda-2a_\lambda+\deg p
\]
by \cite[Lemmas~2.2 and 2.5--2.6]{BFN18}.  These are exactly the degrees
of the numerator, tangent denominator, and dressing in
Equation~\eqref{eq:minuscule}.  The bare matter-free torus class
$e^\lambda$ has degree zero.  This proves the stated grading comparison.
General dressings are decomposed into homogeneous summands.
The additional shifts from a common finite section cutoff will be kept
explicitly in Section~\ref{sec:correspondences}.
No nonnegative grading on $A_F$ itself is assumed.
\end{proof}

Every nontrivial minuscule generator in Proposition~\ref{prop:minuscule} has nonzero
charge.  Consequently the coefficient map $R\to B_F$ is surjective,
and $B_F$ is nonnegatively graded.  More precisely, the ideal in
\eqref{eq:full-fixed-quotient} is generated by those minuscule classes:
after killing them the algebra is generated by $R$ and is supported in
charge zero.  In particular, a coefficient relation in $B_F$ can be
expressed as a finite sum of neutral words of positive length in the
nonzero minuscule generators, with coefficients in $R$.

\begin{lemma}[Flavor base change]\label{lem:flavor-base-change}
Let $F\to F'$ be a homomorphism of algebraic tori, and suppose the action
on $N$ is obtained by restriction of an $F'$-action commuting with $G$.
Then the natural maps are isomorphisms
\begin{align*}
 A_{F'}\otimes_{H_{F'}^*(\mathrm{pt})}H_F^*(\mathrm{pt})&\ \cong A_F,\\
 B_{F'}\otimes_{H_{F'}^*(\mathrm{pt})}H_F^*(\mathrm{pt})&\ \cong B_F.
\end{align*}
They respect the coefficient maps, both gradings, and convolution.
\end{lemma}

\begin{proof}
Filter the affine Grassmannian by bounded unions of orbits and truncate
all section lattices between sufficiently high positive and negative
powers of $t$.  Over an orbit, the truncated space of triples is a
vector bundle of constant relative rank.  The orbit is an affine-space
bundle over a partial flag variety; its equivariant homology, with the
relative shifts in the definition of $A_F$, is even.  Its $F'$-equivariant
homology is free over $H_{F'}^*(\mathrm{pt})$, and restricting the flavor
action gives the displayed change of coefficients on this stratum.
The exact sequences for adjoining strata have no odd terms, so these
statements extend through each finite filtration.  The free successive
quotients also show that tensoring the sequences introduces no Tor term.
Passing to the stabilized truncations and their exhaustive union proves
the first assertion.  Pullback, refined intersection, and proper
pushforward defining convolution commute with this change of
equivariance, so the assertion is an algebra isomorphism.

Flavor parameters have charge zero.  Under the first isomorphism the
ideal generated by all nonzero charge spaces is the extension of the
same ideal before base change.  Taking the quotient therefore gives
the second assertion; no flatness of $B_{F'}$ is required for this step.
\end{proof}

For the moment-map comparisons below, we use the topological affine
Kempf--Ness theorem: the quotient of a character-shifted real moment
level by a compact form of $G$ is homeomorphic to the corresponding
affine-character GIT quotient
\cite[Theorem~4.2(v)]{Hoskins14}.  The same statement restricts to a closed
invariant subvariety.  Indeed, its GIT quotient embeds as a closed
subspace of the ambient quotient, and its orbit closures and their
closed representatives remain in that subvariety.  This applies in
particular to every central complex moment fiber in $T^*N$.

\begin{lemma}[Local constancy of regular stability]\label{lem:local-stability}
Let $\theta$ be a rational character parameter for which every semistable
point of $\mu^{-1}(0)$ is stable and has trivial stabilizer.  There is a
neighborhood $U$ of $\theta$ in $X^*(G)\otimes_{\Z}\mathbb R$ such that
the stable and semistable loci are unchanged for every rational
$\theta'\in U$.  One may choose such a $\theta'$ outside all proper
hyperplanes $\theta'\cdot\mathbf b=0$ with $0<\mathbf b\leq\mathbf v$,
and then clear its denominators by a positive integer.  Empty loci are included.
\end{lemma}

\begin{proof}
We identify a positive multiple of a rational character with the same
stability condition.  Fix a maximal torus of $G$ and the finite set of its
weights on $N\oplus N^*$.  For a point $x$ and $g\in G$, let $C(gx)$ be the
cone generated by the weights occurring in $gx$.  The affine
Hilbert--Mumford criterion says that $x$ is semistable precisely when the
character belongs to every such cone, and is stable with finite stabilizer
precisely when it belongs to all their full interiors in the real weight
space; see \cite{King94} and
\cite[Propositions~2.5 and 2.7]{Hoskins14}.
Although $g$ varies, only finitely many cones
occur.  A sufficiently small neighborhood of $\theta$ remains inside each
cone whose interior contains $\theta$, and outside each cone that does not
contain $\theta$.  A $\theta$-stable point therefore remains stable, and a
$\theta$-unstable point remains unstable.  Regularity accounts for all
points of the closed subvariety $\mu^{-1}(0)$.  This proves the assertion
about the entire loci, including when they are empty.  The geometric
quotient and its principal gauge bundle are unchanged, as is freeness.
Choose a rational point in this neighborhood avoiding the finitely many
proper hyperplanes specified in the statement, and then rescale by a
positive integer.  If $G$ is trivial, its character space is a point and
the list of bounded nonzero dimension vectors is empty, so the assertion
is immediate.
\end{proof}

The Kirwan map $R\to H_F^*(X)$ is surjective.  For nonempty $X$,
Lemma~\ref{lem:local-stability} permits a rational perturbation $\theta$
of $\nu$ that preserves the entire stable locus $D$ and avoids
$\theta\cdot\mathbf b=0$ for $0<\mathbf b\leq\mathbf v$.  Its extension
$(\theta,-\theta\cdot\mathbf v)$ is nondegenerate for the framed
dimension vector $(\mathbf v,1)$: a proper subvector has final
coordinate zero, or its nonzero complement does.  Clearing denominators
preserves stability.  The quotient, its principal gauge bundle, and its
tautological bundles are unchanged, so ordinary Kirwan surjectivity
applies with the same coefficient map
\cite[Definition~3.1 and Theorem~1.2]{MN18}; see also
\cite[Section~2.1]{DK25}.  The torus-equivariant conclusion is
\cite[Corollary~1.5]{MN18}; we recall its argument for an arbitrary
commuting flavor torus.  Choose a homogeneous basis of $H^*(X)$ represented by polynomials
in the ordinary tautological Chern classes.  These classes lift to
$F$-equivariant Chern classes.  Leray--Hirsch for
$X\to EF\times_F X\to BF$ therefore makes $H_F^*(X)$ free over
$H_F^*(\mathrm{pt})$ on these lifts and proves the asserted equivariant
surjectivity.  In particular $H^*(X)$ is finite-dimensional and even.
For empty $X$ the assertion means the unique surjection to the zero ring.

\section{A finite correspondence test for Coulomb relations}
\label{sec:correspondences}

We work with the finite-dimensional algebraic group $\mathsf K=G\times F$.
All homology in this section is equivariant Borel--Moore homology, with
complex coefficients.  We use constructible sheaves only on finite-dimensional
varieties.  The purpose of the section is to express a coefficient relation
in the Coulomb algebra as a relation between such sheaves.

\subsection{Bounded lattice families}

A point $b=gG(\OO)$ of $\Gr_G$ determines the lattice $gN(\OO)$ in
$N(\KK)$.  For a BFN triple $[g,s_0]$, we use $n=gs_0$ as section
coordinate; the fiber over $b$ is $N(\OO)\cap gN(\OO)$.
Let $B$ be a projective $\mathsf K$-variety equipped with a
$\mathsf K$-equivariant map to $\Gr_G$ with bounded image.  The map need not
be an embedding.  For any finite collection of these families, choose $l$
large enough that all their lattices contain $t^lN(\OO)$ and are contained
in $t^{-l}N(\OO)$.  Set
\begin{equation}
 \label{eq:finite-lattice-space}
 \begin{aligned}
 L&=t^{-l}N(\OO)/t^lN(\OO),&
 E_b&=gN(\OO)/t^lN(\OO),\\
 S_B&=\{(b,n)\in B\times L:n\in E_b\},&
 f_B&:S_B\longrightarrow L.
 \end{aligned}
\end{equation}
Here and below the same notation is used for a lattice and for its pullback
to a parameter space.  The representation map
$\Gr_G\to\Gr_{\GL(N)}$ identifies $S_B$ with the pullback of a
tautological vector bundle on a finite Grassmannian.  We work separately on
components if its rank varies.  If $\chi_N=\det(N|_G)$, its rank on a
component of degree $\gamma\in\pi_1(G)$ is
\begin{equation}
 \label{eq:lattice-rank}
 r_B=l\dim N-\langle\chi_N,\gamma\rangle.
\end{equation}
The incidence $S_B\subset B\times L$ is closed, so $f_B$ is proper.
We write $S_{\mathrm{pt}}=N(\OO)/t^lN(\OO)$ for the family at the
identity lattice and $i=f_{\mathrm{pt}}$.  Put $r=l\dim N$ and
$\mathcal F_B=Rf_{B,*}\C_{S_B}$; in particular,
$\mathcal F_{\mathrm{pt}}=i_*\C_{S_{\mathrm{pt}}}$.

Let $Z$ be a finite union of affine Schubert varieties in the neutral
component, containing the identity lattice.  Restriction to that lattice
gives a degree-zero morphism
\begin{equation}
 \label{eq:restriction-at-identity}
 \rho_Z:\mathcal F_Z\longrightarrow\mathcal F_{\mathrm{pt}}.
\end{equation}
Both section bundles here have rank $r$.  The base $Z$ is allowed to be
singular.

The notation $\Hom^d_{\mathsf K,L}(\mathcal F,\mathcal G)$ denotes
$\Hom(\mathcal F,\mathcal G[d])$ for $\mathsf K$-equivariant constructible
complexes on $L$, computed on the common Borel approximations used below.

\begin{lemma}[Finite correspondence factorization]
 \label{lem:correspondence-factorization}
Let $u\in R$ be homogeneous of cohomological degree $d$, and suppose that
in $A_F$ it has an expression
\begin{equation}
 \label{eq:neutral-word-relation}
 u=\sum_w M_{\lambda_{w,1},p_{w,1}}\cdots
                 M_{\lambda_{w,k_w},p_{w,k_w}},
 \qquad k_w\geq1,
 \qquad \sum_j[\lambda_{w,j}]=0\in\pi_1(G).
\end{equation}
Every factor is a nontrivial dressed positive or negative minuscule
generator of Proposition~\ref{prop:minuscule}, and every word is homogeneous
of degree $d$.  Coefficients in $R$ may be included in the dressings.
Choose an ordering of the factors of each word.  There are $Z$ as above
bounding their final lattices and a common truncation
\eqref{eq:finite-lattice-space} such that
\begin{equation}
 \label{eq:finite-relation-factorization}
 u\rho_Z=\sum_w\Psi_w
 \quad\text{in}\quad
 \Hom^d_{\mathsf K,L}
       (\mathcal F_Z,\mathcal F_{\mathrm{pt}}).
\end{equation}
Each $\Psi_w$ factors through a cohomological shift of
$\mathcal F_{\Gr_G^{\lambda_{w,1}}}$.
Only the endpoints are required to lie in $Z$ or in the neutral component;
all intermediate lattices are included in the common truncation.
\end{lemma}

\begin{proof}
We first specify the correspondences and their degrees, then prove the
identity by a faithful localization calculation.

\paragraph{Correspondences and shifts.}
For proper maps $f:X\to L$ and $g:Y\to L$, write $\omega_Y$ for the
complex-linear dualizing complex of $Y$.  Proper base change and adjunction
give
\[
 R\Hom_L(Rf_*\C_X,Rg_*\omega_Y)
 \simeq R\Gamma(X\times_LY,\omega_{X\times_LY}).
\]
For example, if $q:X\times_LY\to Y$ is the projection, the left side is
$R\Hom_Y(q_*\C_{X\times_LY},\omega_Y)$, which gives the displayed
identity by $q_*\dashv q^!$.  Thus, when $Y$ is smooth of complex dimension
$e$, we have
\begin{equation}
 \label{eq:finite-correspondence-Hom}
 \Hom^a_{\mathsf K,L}(Rf_*\C_X,Rg_*\C_Y)
 \simeq H^{\mathsf K}_{2e-a}(X\times_LY).
\end{equation}
The equivariant version follows on common Borel approximations; their
additional dimensions cancel in the shift.  These finite-type models have
triangulable analytic topology and finite cohomological dimension.
We use the proper-support base-change and adjunction formalism of
\cite[Theorems~5.10 and 7.7]{SchnurerSoergel16}, with complex coefficients,
and the duality and Borel--Moore identifications of
\cite[Example~2.7(3) and Corollary~2.11]{MaximSchurmann21}.
In particular, the source
$X$ need not be smooth.  Composition across a smooth middle space is
refined intersection along its diagonal followed by proper pushforward.
These operations are the adjunction and base-change operations in the
displayed identification.

Fix one word, temporarily suppressing its index $w$.  Let $B_j$ be the
variety of chains of lattices with successive relative positions
$\lambda_1,\ldots,\lambda_j$, starting at the identity.  Set
$B_0=\mathrm{pt}$.  Each forgetful map $B_j\to B_{j-1}$ is a bundle with
fiber the closed minuscule orbit $\Gr_G^{\lambda_j}$, so the $B_j$ are
smooth and projective.  Write
\[
 a_j=\dim\Gr_G^{\lambda_j},\qquad
 r_j=\rk E_{b_j},\qquad
 I_j=\{(b_0,\ldots,b_j,n):n\in E_{b_{j-1}}\cap E_{b_j}\}.
\]
The intersection has constant rank by fixed relative position.  It defines
a vector bundle on $B_j$ and a correspondence
\[
 S_{B_j}\ \overset{v_j}{\hookleftarrow}\ I_j
       \xrightarrow{u_j} S_{B_{j-1}}
\]
over $L$, where $v_j$ is the closed subbundle inclusion and $u_j$ forgets
the last lattice.  Write
$c_j=\rk(E_{b_j}/I_j)$, using the same symbol $I_j$ for the intersection
bundle.  The dressing $p_j$ gives an equivariant characteristic class on
$B_j$: use the two tautological blocks of the minuscule modification in
the preceding lattice frame, together with the flavor classes.  This is
the twisted-orbit version of the stabilizer-invariant dressing in
\eqref{eq:minuscule}.

If $\delta_j=\deg p_j$, the class $p_j\cap[I_j]$ defines a morphism
$\mathcal F_{B_j}\to\mathcal F_{B_{j-1}}[q_j]$, where
\begin{equation}
 \label{eq:chain-shift}
 q_j=2(r_{j-1}-r_j+c_j-a_j)+\delta_j.
\end{equation}
Indeed, $\dim I_j=\dim B_{j-1}+a_j+r_j-c_j$, and
\eqref{eq:finite-correspondence-Hom} uses the target dimension
$\dim S_{B_{j-1}}=\dim B_{j-1}+r_{j-1}$.
The endpoint map $B_k\to Z$ induces a proper map $S_{B_k}\to S_Z$.
Its pullback morphism on constant sheaves, followed by pushforward to
$L$, gives $\mathcal F_Z\to\mathcal F_{B_k}$ with no shift.
Compose it with the $k$ incidence morphisms to obtain $\Psi_w$.
This already gives the stated first-step factorization.

At a modification of type $\lambda_j$,
$c_j=\sum_\beta(-\langle\beta,\lambda_j\rangle)_+$, with multiplicities
and with the pairing taken on the gauge part.  The rational coefficient
in Equation~\eqref{eq:minuscule} has degree $2c_j-2a_j+\delta_j$,
counting linear forms with degree two.  For the neutral word, the sum of
these degrees is its BFN degree $d$, by
Proposition~\ref{prop:minuscule}.  Moreover,
Equation~\eqref{eq:lattice-rank} gives $r_k=r_0=r$.  Consequently
\[
 \sum_jq_j
 =2(r_0-r_k)+\sum_j(2c_j-2a_j+\delta_j)=d.
\]
Thus $\Psi_w$ and $u\rho_Z$ belong to the same graded Hom space.
No equality of intermediate ranks or identification of nonneutral degree
conventions is used.

\paragraph{Faithfulness of the localization test.}
Put $C_Z=S_{\mathrm{pt}}\times_LS_Z$ and let $\iota:C_Z\hookrightarrow
S_Z$ be the inclusion.  The identity
\eqref{eq:finite-correspondence-Hom}, with smooth target
$S_{\mathrm{pt}}$, identifies the Hom space in
\eqref{eq:finite-relation-factorization} with
$H^{\mathsf K}_{2r-d}(C_Z)$.  Under it, $\rho_Z$ is the fundamental class
of $S_{\mathrm{pt}}$ over the identity lattice.

Let $P=H^*_{T\times F}(\mathrm{pt})$.  The stratification of $C_Z$ over
the finitely many $G(\OO)$-orbits in $Z$ has vector-bundle strata.  Each
orbit is an affine-space bundle over a partial flag variety and contracts
to it under loop scaling $t\mapsto at$ as $a\to0$.  This contraction
commutes with $T$, and $F$ acts trivially on the orbit.  Its ordinary
equivariant cohomology is therefore even and free over $P$; smooth
equivariant Poincar\'e duality gives the same assertion for Borel--Moore
homology, with a dimension shift.  The section bundle adds an even Thom
shift.
The closed-stratum exact sequences split into short exact sequences by
parity, and their free quotients make each extension free over $P$.
It follows that $H_*^{T\times F}(C_Z)$ is free, in particular
torsion-free.  This is the finite version of the orbit-filtration argument
in \cite[Lemmas~2.5--2.6]{BFN18}.

The $T$-fixed lattices in $\Gr_G$ are precisely the diagonal lattices
$t^\eta$: a fixed lattice in each $V_i(\KK)$ splits along its distinct
diagonal characters, and a lattice in a one-dimensional Laurent series
space is $t^m\OO$.  At $t^\eta$, every $(T\times F)$-weight-zero
summand of $N$ has gauge pairing zero with $\eta$.  Its lattice is
unchanged by intersection with $N(\OO)$.  Thus $\iota$ identifies the
fixed loci of $C_Z$ and $S_Z$.  Equivariant localization makes
$\iota_*$ an isomorphism after tensoring with $\operatorname{Frac}P$;
torsion-freeness makes this localized pushforward a faithful test on
$H_*^{T\times F}(C_Z)$.  Passing to Weyl invariants recovers
$\mathsf K$-equivariance.  This is also the fixed-fiber calculation of
\cite[Lemma~5.13]{BFN18}.

After localization, the classes of the entire section fibers
$E_{t^\eta}$ form a basis for the homology of $S_Z$, indexed by its
fixed lattices.  These classes, rather than isolated points in a fiber,
allow weight-zero summands of $N$ without inverting a zero Euler class.
We denote this basis formally by $e^\eta$.

\paragraph{The chain calculation.}
For each partial chain put $C_j=S_{\mathrm{pt}}\times_LS_{B_j}$, and let
$h_j:C_j\hookrightarrow S_{B_j}$ forget the initial section constraint.
Let $z_j$ be the Borel--Moore class of the composite
$\mathcal F_{B_j}\to\mathcal F_{\mathrm{pt}}[\sum_{h=1}^j q_h]$
formed by the first $j$ incidence morphisms.  Since $S_{B_j}$ is smooth, write
$(h_j)_*z_j=\alpha_j\cap[S_{B_j}]$.  We claim that
\begin{equation}
 \label{eq:chain-euler-product}
 \alpha_k=\prod_{j=1}^k p_j\,\eu(E_{b_j}/I_j),
\end{equation}
with all factors pulled to $B_k$.

The initial class has $\alpha_0=1$.  At the $j$th step, composition is
refined intersection along the diagonal of the smooth middle space
$S_{B_{j-1}}$.  This intersection commutes with the proper pushforward
that forgets the initial section constraint.  We may therefore first
replace $z_{j-1}$ by $(h_{j-1})_*z_{j-1}
=\alpha_{j-1}\cap[S_{B_{j-1}}]$.  The remaining intersection pulls this
cohomology class back to $I_j$, giving
\[
 (h_j)_*z_j
 =(v_j)_*\bigl(u_j^*\alpha_{j-1}\,p_j\cap[I_j]\bigr).
\]
Here the regular diagonal belongs to the smooth middle space; no
smoothness of $u_j$ is needed.  The classes $\alpha_{j-1}$ are pulled
from the chain bases, and the subbundle inclusion $v_j$ multiplies by
$\eu(E_{b_j}/I_j)$.  This proves \eqref{eq:chain-euler-product} by
induction.

Under the same Hom--Borel--Moore identification, precomposition with the
unit for the proper endpoint map $S_{B_k}\to S_Z$ is proper pushforward
along $C_k\to C_Z$.  Forgetting the initial section constraint commutes
with this pushforward.  Thus the endpoint pushforward of
$\alpha_k\cap[S_{B_k}]$ is exactly the image of $\Psi_w$ in the
localization test.

Fixed chains are sequences of Weyl translates $\lambda'_j$ of the chosen
minuscule coweights; their successive lattices are
$t^{\eta_j}$ with $\eta_j=\sum_{h\leq j}\lambda'_h$.
At such a chain the $j$th modification is translated by
$t^{\eta_{j-1}}$.  This translation commutes with the diagonal torus.
There is no loop rotation, so it changes neither the relevant weights nor
the dressing values.  Hence its quotient Euler class is
\[
 \prod_\beta\beta^{(-\langle\beta,\lambda'_j\rangle)_+},
\]
and its relative tangent Euler class is
$\eu_T(T_{t^{\lambda'_j}}\Gr_G^{\lambda_j})$.
Localization on the smooth chain base therefore sends
\eqref{eq:chain-euler-product}, after endpoint pushforward, to
\[
 \sum_{(\lambda'_1,\ldots,\lambda'_k)}
 e^{\lambda'_1+\cdots+\lambda'_k}
 \prod_{j=1}^k
 \left(
 p_{j,\lambda'_j}
 \frac{\prod_\beta\beta^{(-\langle\beta,\lambda'_j\rangle)_+}}
 {\eu_T(T_{t^{\lambda'_j}}\Gr_G^{\lambda_j})}
 \right).
\]
This is exactly the product of the expressions
\eqref{eq:minuscule} for the word.  All dressings are transported with
their coweights.  The fiber ranks at the endpoint equal $r$, so no
additional fiber factor enters endpoint pushforward.

The localized pushforward of $u\rho_Z$ is $u e^0$.
Equation~\eqref{eq:neutral-word-relation} and the abelianization of
Proposition~\ref{prop:minuscule} identify it with the sum of the displayed
word expressions.  The faithful test just proved implies the equality
in $H^{\mathsf K}_{2r-d}(C_Z)$ and hence
\eqref{eq:finite-relation-factorization}.
\end{proof}

\begin{remark}
Every $u\in\ker(R\to B_F)$ admits an expression of the form
\eqref{eq:neutral-word-relation}, degree by degree.  Indeed,
Proposition~\ref{prop:minuscule} implies that the ideal of all nonzero
topological degrees is generated by the nontrivial positive and negative
minuscule generators.  Expand a relation in that ideal and retain its
neutral component.  Commutativity of $A_F$ permits any desired ordering
of the factors before applying Lemma~\ref{lem:correspondence-factorization}.
\end{remark}

The relation has now been transferred to a finite-dimensional category,
with the identity-lattice map on one side and first-step factorizations
on the other.  The stable-pair test in Section~\ref{sec:stable-pairs}
will annihilate the selected first-step objects while keeping a split
restriction from the identity-lattice test onto cohomology.

\section{Testing coefficient relations by stable pairs}
\label{sec:stable-pairs}

We use the finite-dimensional correspondences of
Section~\ref{sec:correspondences} to prove
\begin{proposition}
\label{prop:kernel-inclusion}
For every datum and regular stability in Theorem~\ref{thm:main}, the
coefficient map induces a graded $R$-algebra surjection
\begin{equation}
\label{eq:stable-pair-surjection}
 B_F\twoheadrightarrow H_F^*(X).
\end{equation}
\end{proposition}
Throughout this section, $\mathsf K=G\times F$.  We retain the notation
$L$, $E_b$, $S_B$, $f_B$, $i=f_{\mathrm{pt}}$, and $\rho_Z$ from
Lemma~\ref{lem:correspondence-factorization}.  All cohomology in the test
below is ordinary equivariant cohomology.

\subsection{A degree constraint on global stable pairs}

We first record the geometric consequence of stability that will exclude
the first lattice of each word.  For a principal $G$-bundle $\mathcal P$
on $\mathbb P^1$, write $\mathcal L_\nu$ for its character line bundle,
with the convention that evaluating a $\nu^m$-semi-invariant on a section
of an associated representation gives a section of
$\mathcal L_\nu^{\otimes m}$.  Its degree is a linear function of the
topological class of $\mathcal P$.

\begin{lemma}
\label{lem:stable-section-degree}
Let $\mathcal P$ be a principal $G$-bundle on $\mathbb P^1$, equipped with
a frame at infinity, and let $\sigma$ be a global section of
$\mathcal P\times_G(T^*N)$.  Suppose that its value $\sigma_\infty$ in this
frame is $\nu$-stable and has trivial stabilizer.  Then
$\deg\mathcal L_\nu\geq0$.  If $\deg\mathcal L_\nu=0$, the pair
$(\mathcal P,\sigma)$ has a unique trivialization, agreeing with the given
frame at infinity, in which $\sigma$ is the constant section
$\sigma_\infty$.

In particular, suppose that $\mathcal P$ is specified by a lattice
$b\in\Gr_G$, with its usual trivialization on $\mathbb P^1\setminus\{0\}$.
In the degree-zero case its topological class is zero and
\begin{equation}
\label{eq:based-loop-trivialization}
 b=hG(\OO),\qquad
 h\in G(\C[t^{-1}]),\qquad h(\infty)=1.
\end{equation}
In the off-zero frame, the section is $h\sigma_\infty$.
\end{lemma}

\begin{proof}
Choose a $\nu^m$-semi-invariant $f$, with $m>0$, that is nonzero at
$\sigma_\infty$.
Its evaluation on $\sigma$ is a global section of a positive power of
$\mathcal L_\nu$, nonzero at infinity.  This proves nonnegativity of the
degree.

If the degree is zero, this section is nowhere zero.  Thus $\sigma$ is
everywhere semistable and determines a morphism
\[
 \gamma:\mathbb P^1\longrightarrow (T^*N)/\!/_{\nu}G.
\]
Every invariant polynomial evaluated on $\sigma$ is a regular function
on $\mathbb P^1$, hence constant.  The image of $\gamma$ therefore lies
in one fiber of the projective morphism to the affine quotient
$(T^*N)/\!/G$.  A sufficiently divisible GIT polarization is relatively
ample and pulls back to a positive power of $\mathcal L_\nu$.  Its
degree is zero, so $\gamma$ is constant: an ample line bundle has
positive degree on the image of a nonconstant morphism from
$\mathbb P^1$.  These are the usual character-linearized affine GIT
constructions; see \cite[Section~2]{King94}.

The constant value is the stable quotient point represented by
$\sigma_\infty$.  Its fiber in the semistable locus is the single orbit
$G\sigma_\infty$.  Consequently $\sigma$ is a section of the associated
orbit bundle $\mathcal P\times_G(G\sigma_\infty)$.  Since
$G\sigma_\infty\simeq G$ equivariantly, this section trivializes
$\mathcal P$.  Freeness makes the trivialization unique once the
coordinates of the section are prescribed to be $\sigma_\infty$; it
agrees with the specified frame at infinity.

In the lattice presentation, compare this global trivialization with the
fixed trivialization off zero.  Their transition is a regular map
$\mathbb A^1_{t^{-1}}\to G$ equal to $1$ at infinity.  It is the loop $h$
in \eqref{eq:based-loop-trivialization}, and the lattice of regular
sections at zero is $hG(\OO)$.
\end{proof}

Every point of $D=\mu^{-1}(0)^{\nu\text{-stable}}$ is stable with trivial
stabilizer also in the ambient representation $T^*N$.  Indeed, for a
closed invariant subvariety, semi-invariants lift by exactness of
reductive invariants, and orbit closures stay in that subvariety.
The affine GIT stability criterion therefore agrees with its ambient
restriction.  Lemma~\ref{lem:stable-section-degree} applies to the
\emph{pair} $(s,p)\in D$.

\subsection{The finite vanishing-cycle test}

For the truncation parameter $l$ of
Lemma~\ref{lem:correspondence-factorization}, put
\[
 L_-=\bigoplus_{j=-l}^{-1}t^jN,\qquad
 L_+=\bigoplus_{j=0}^{l-1}t^jN,\qquad
 P_-=\bigoplus_{j=-l}^{-1}t^jN^*.
\]
Thus $L=L_-\oplus L_+$, and the residue pairing identifies
$P_-$ with $L_+^*$.  Regard $s\in N$ and $p\in N^*$ as constant series.
On $\widetilde L=L\times D\times P_-$ define the $\mathsf K$-invariant
function
\begin{equation}
\label{eq:stable-test-function}
 W(n,s,p,y)=
 \Res\bigl(\langle p+y,n\rangle-\langle y,s\rangle\bigr),
 \qquad (s,p)\in D.
\end{equation}
After restricting $n$ to a moving lattice, variation in $y$ imposes the
prescribed nonnegative modes $\operatorname{pr}_{L_+}n=s$, while variation
in $n$ imposes dual-lattice membership for $p+y$.  These are the extension
conditions for the two components of the stable pair across the lattice
modification.
The residue is the coefficient of $t^{-1}$; the expression is
well-defined on the truncated spaces.

Use the shifted vanishing-cycle functor $\Phi_f=\phi_f[-1]$, normalized
so that $\Phi_0$ is the identity.  Vanishing cycles will be regarded
as complexes on the ambient space by closed direct image from the zero
fiber.  The test of a complex $\mathcal F$ on $L$ is
\begin{equation}
\label{eq:stable-test-functor}
 \mathsf T(\mathcal F)
 =R\Gamma_{\mathsf K}\!
 \left(\widetilde L,\,
 \Phi_W\,\operatorname{pr}_L^*\mathcal F\right).
\end{equation}
For a bounded projective lattice family $B$, set
$M_B=S_B\times D\times P_-$, and let $W_B$ be the pullback of $W$.
Proper base change and compatibility of vanishing cycles with proper
direct image give
\begin{equation}
\label{eq:stable-test-proper}
 \mathsf T(f_{B,*}\C_{S_B})
 \simeq R\Gamma_{\mathsf K}
       (M_B,\Phi_{W_B}\C_{M_B}).
\end{equation}
Here proper-direct-image compatibility is
\cite[Proposition~4.19(1)]{MaximSchurmann21}; the smooth projection
pullback uses part~(2).  The same operations in the critical-cohomology
setting are recalled in \cite[Appendix~A.1]{BottaTamagni26}.

The test is $R$-linear.  To see this precisely, work on a finite Borel
approximation with base $B\mathsf K$.  A coefficient class $u$ is a
morphism $\C_{B\mathsf K}\to\C_{B\mathsf K}[\deg u]$, pulled back to each
space.  The projection formula makes proper direct image compatible
with this action \cite[Theorem~6.3]{SchnurerSoergel16}.
Nearby cycles commute naturally with tensoring by a
locally constant complex pulled back from the base; taking the
specialization cone gives the same assertion for vanishing cycles.
Naturality with respect to the morphism representing $u$ therefore
identifies its tested action with cup product by the restricted class.
Passing to sufficiently large approximations proves the assertion for
$R=H_{\mathsf K}^*(\mathrm{pt})$.

\begin{lemma}
\label{lem:stable-test-support}
The support of $\Phi_{W_B}\C_{M_B}$ is contained in the closed locus
$\Sigma_B$ defined by
\begin{equation}
\label{eq:stable-critical-equations}
 n\in s+L_-,
 \qquad
 p+y\in gN^*(\OO)/t^lN^*(\OO),
 \qquad b=gG(\OO).
\end{equation}
Here $n\in E_b$ is already part of the definition of $M_B$.
If $B$ lies in a nonzero topological component whose character-line
degree is nonpositive, then $\Sigma_B$ is empty and the test in
\eqref{eq:stable-test-proper} is zero.
\end{lemma}

\begin{proof}
Differentiating in $y$ gives
$\operatorname{pr}_{L_+}n=s$.  Differentiating in the section variable
$n\in E_b$ says that $p+y$ annihilates $E_b$ in the residue pairing.
The annihilator is the displayed truncated dual lattice.  These
conditions are zeros of algebraic bundle maps and hence define a closed
locus.

This derivative calculation gives a support bound even when $B$ is
singular.  Off $\Sigma_B$, a derivative in one of the vector-bundle fiber
coordinates is nonzero.  The function is affine linear in that
coordinate; where its coefficient is invertible, replacing the
coordinate by $W_B$ gives a product chart over the unchanged base.
The constant sheaf has zero vanishing cycles in this chart.

A bounded lattice in each general linear factor specifies a modification
of the trivial vector bundle at zero, determined by finite jets.
A point of $\Sigma_B$ gives two global sections on $\mathbb P^1$ for the
resulting principal bundle.  In the off-zero frame they are the
polynomials $n$ and $p+y$ in $t^{-1}$, with values $s$ and $p$ at
infinity.  The lattice conditions are exactly their extendability at
zero: the discarded $t^l$ tails already lie in the corresponding
lattices.  Their value at infinity is the stable pair $(s,p)$.
Lemma~\ref{lem:stable-section-degree} now excludes negative
character-line degree and, in degree zero, forces the entire
topological class to be zero.  This proves the last assertion.
\end{proof}

\subsection{The neutral component and its quadratic normal form}

Let $Z$ be the neutral projective Schubert union used in
Lemma~\ref{lem:correspondence-factorization}, containing the base lattice
$1$.  Set $r=l\dim N$.  Since the index of the associated representation
lattice factors through $\pi_1(G)$, the bundles $E_b$ have rank $r$ on
$Z$.  Define the invariant open subset
\begin{equation}
\label{eq:stable-test-open}
 U=\{b\in Z:\operatorname{pr}_{L_+}:E_b\xrightarrow{\sim}L_+\}.
\end{equation}
It contains $1$.

\begin{lemma}
\label{lem:quadratic-test}
Assume $D\neq\varnothing$.  The locus $\Sigma_Z$ is a closed graph
$\Gamma\simeq U\times D$ in $M_Z$.  There are compatible
$\mathsf K$-equivariant identifications
\begin{equation}
\label{eq:quadratic-vanishing-cycles}
 \Phi_{W_Z}\C_{M_Z}
 \simeq \iota_{\Gamma,*}\C_\Gamma[-2r],
 \qquad
 \Phi_{W_{\mathrm{pt}}}\C_{M_{\mathrm{pt}}}
 \simeq \iota_{D,*}\C_D[-2r].
\end{equation}
Under these identifications, the test of $\rho_Z$ induces ordinary
restriction
\begin{equation}
\label{eq:quadratic-restriction}
 H_{\mathsf K}^*(U\times D)\longrightarrow H_{\mathsf K}^*(D)
\end{equation}
with the same uniform shift on source and target.  This map is split
surjective and $R$-linear.
\end{lemma}

\begin{proof}
First, every point of $\Sigma_Z$ lies over $U$.  By
Lemma~\ref{lem:stable-section-degree}, its principal bundle is trivial,
with a based loop $h$ as in \eqref{eq:based-loop-trivialization}.
A vector in the kernel of $E_b\to L_+$ is a negative polynomial section
of its associated $N$-bundle, extending at zero and vanishing at
infinity.  That associated bundle is trivial, so the vector is zero.
The two spaces have dimension $r$, and the projection is an
isomorphism.

Over $U$, write $E_b$ as the graph of a regular linear map
$A_b:L_+\to L_-$.  Thus $n=a+A_ba$ for $a\in L_+$.  Define the linear
functional
\[
 \ell_{b,p}(a)=\Res\langle p,A_ba\rangle.
\]
Via $P_-=L_+^*$, Equation~\eqref{eq:stable-test-function} becomes
\begin{equation}
\label{eq:stable-quadratic-coordinates}
 W_Z=\langle y,a-s\rangle+\ell_{b,p}(a).
\end{equation}
Its fiber critical equations have the unique solution
$a=s$, $y=-\ell_{b,p}$, depending regularly on $(b,s,p)$.
Consequently $\Sigma_Z$ is the graph $\Gamma\simeq U\times D$.
It is closed in all of $M_Z$ because the equations defining $\Sigma_Z$ are
closed and every solution has already been shown to lie over $U$.

It remains to compute the critical value.  Apply
Lemma~\ref{lem:stable-section-degree} to this solution.  In the
off-zero frame,
\[
 n=hs,\qquad p+y=hp.
\]
Invariance of the pairing gives
$\langle hp,hs\rangle=\langle p,s\rangle$, a constant series.  Hence
\[
 W_{\mathrm{crit}}
 =-\langle[t^{-1}](hp),s\rangle.
\]
Write $h=1+\xi t^{-1}+O(t^{-2})$, where $\xi\in\mathfrak g$.
The coefficient $[t^{-1}](hp)$ is the infinitesimal dual action
$\xi\cdot p$.  Therefore
\begin{equation}
\label{eq:stable-critical-value}
 W_{\mathrm{crit}}
 =-\langle\xi\cdot p,s\rangle
 =\langle p,\xi\cdot s\rangle
 =\langle\mu(s,p),\xi\rangle=0.
\end{equation}
Thus the regular, equivariant translations
$a'=a-s$, $y'=y+\ell_{b,p}$ turn
\eqref{eq:stable-quadratic-coordinates} into the fixed split form
\[
 W_Z=\langle y',a'\rangle
\]
over $U\times D$.

The vanishing cycles of this form on $L_+\oplus L_+^*$ are one copy of
$\C[-2r]$ at the origin.  The external-product theorem applies to the
constant sheaf of $U\times D$ even if $U$ is singular
\cite[Theorem, pp.~1--2 of the arXiv preprint]{Massey99}.  The splitting is global and
equivariant.  Its structure group acts as $g\oplus g^{-T}$, so the
quadratic vanishing-cycle line has trivial equivariant monodromy.
Choosing its identification with $\C$ once gives compatible
identifications over the whole base.

Lemma~\ref{lem:stable-test-support} says that the vanishing cycles are
zero off the closed graph $\Gamma$.  The product calculation therefore
gives the first isomorphism in
\eqref{eq:quadratic-vanishing-cycles} on all of $M_Z$.  In particular,
there are no further cycles over $Z\setminus U$: every point there
fails at least one of the fiber equations.  The same computation at
$1$ has $A_1=0$, $\ell_{1,p}=0$, and the same $2r$ quadratic variables,
giving the second isomorphism.

The map $\rho_Z$ is restriction to the base lattice.  To identify its
test as a morphism, not just its source and target, put $B=U\times D$
and consider the maps
\[
 D\xrightarrow{\,j\,}U\times D\xrightarrow{\,\pi\,}D,
 \qquad j(d)=(1,d),\quad \pi(b,d)=d.
\]
These maps are $\mathsf K$-equivariant and satisfy $\pi j=\id$.  Their ordinary
cohomological pullbacks satisfy $j^*\pi^*=\id$.
On the same base $B$, apply the natural external-product identification
to the morphism $\C_B\to Rj_*\C_D$, with zero function on $B$ and the
fixed split quadratic on $L_+\oplus L_+^*$.
Since $j$ is closed, proper-direct-image compatibility identifies the
target with the quadratic calculation over $D$.  One common choice of
generator of the quadratic line then identifies the resulting morphism
with $\C_B[-2r]\to Rj_*\C_D[-2r]$.
Taking ordinary equivariant cohomology gives exactly
Equation~\eqref{eq:quadratic-restriction}.  This uses naturality in a
sheaf morphism and proper direct image, not commutation with an arbitrary
nonsmooth pullback.  The identity $j^*\pi^*=\id$ supplies the
$R$-linear splitting, without any properness requirement on $\pi$.
\end{proof}

\begin{proof}[Proof of Proposition~\ref{prop:kernel-inclusion}]
The assertion is immediate if $D=\varnothing$, so assume $D$ is
nonempty.  Let $u\in R$ be homogeneous and zero in $B_F$.
Proposition~\ref{prop:minuscule} expresses $u$ in $A_F$ as a finite sum
of positive-length words in nontrivial dressed positive and negative
minuscule generators, with coefficients absorbed into the dressings.
Taking the full topological-degree-zero part, every word has total
topological degree zero.

The character-line degrees of its steps sum to zero.  Choose a step
with nonpositive character-line degree and, by commutativity, put it
first.  Its topological degree is nonzero.  Apply
Lemma~\ref{lem:correspondence-factorization}, with a common neutral
endpoint bound $Z$.  The morphism $u\rho_Z$ is a sum of compositions,
each factoring through the first-step object
$f_{B,*}\C_{S_B}$, with the prescribed shifts.  By
Lemma~\ref{lem:stable-test-support}, the test of every such object is
zero.  It follows that
\[
 \mathsf T(u\rho_Z)=0.
\]

The degree calculation in
Lemma~\ref{lem:correspondence-factorization} identifies the degree of
this morphism with $\deg u$, including when intermediate lattice
components are nonneutral.  The test commutes with shifts, and
Lemma~\ref{lem:quadratic-test} places the same quadratic shift on its
source and target.  Its $R$-linearity therefore gives the equality of
graded maps
\[
 (u|_D)\smile j^*=0.
\]
Apply this equality to $\pi^*(1)$.  Since $j^*\pi^*(1)=1$, we obtain
$u|_D=0$ in $H_{\mathsf K}^*(D)$.

The free stable-locus action identifies
$H_{\mathsf K}^*(D)=H_F^*(X)$, compatibly with the coefficient map.
We have proved
\[
 \ker(R\to B_F)\subseteq\ker(R\to H_F^*(X)).
\]
Both rings are generated by their coefficient images.  For $B_F$ this
follows from Proposition~\ref{prop:minuscule}; for $H_F^*(X)$ it follows
from Kirwan surjectivity \cite{MN18} and its equivariant extension
recalled above.  The displayed inclusion induces the graded $R$-algebra
surjection
\eqref{eq:stable-pair-surjection}.
\end{proof}

\section{A Laurent model and a filtration by multipartitions}
\label{sec:laurent}

To complement the surjection of Proposition~\ref{prop:kernel-inclusion},
we seek an upper bound for the Hilbert series of $B_0$.  One loop at each
vertex cancels the minuscule tangent denominators, giving a Laurent model
for this calculation.

Throughout this Section the flavor torus is trivial, and the quiver has at
least one loop at every vertex of positive gauge dimension.  Delete the
vertices of gauge dimension zero, and choose one loop at every remaining
vertex.  Let $Q'$ be the quiver obtained by deleting these chosen loops.
Write
\[
 n_{ij}=\#Q'_1(i,j),\qquad c_{ij}=n_{ij}+n_{ji},\qquad
 C(\mathbf v)=\frac12\sum_{i,j}c_{ij}v_iv_j.
\]
In particular, $c_{ii}=2n_{ii}$.  We use the set
\[
 \Omega=\{(i,a):i\in Q_0,\ 1\leq a\leq v_i\}
\]
of gauge slots, the polynomial ring $P=\C[x_{ia}:(i,a)\in\Omega]$, and
$\mathcal W=\prod_iS_{v_i}$.  Let $\mathcal B'$ be the multiset consisting
of the weights $x_{jb}-x_{ia}$ of all arrows of $Q'$ and $w_i$ copies of
$-x_{ia}$ at each gauge slot.  The zero weights arising from diagonal
entries of loops can be omitted: their coweight pairing is always zero.

For $\lambda\in\Z^\Omega$, put
\[
 E_\beta(\lambda)=\bigl(-\langle\beta,\lambda\rangle\bigr)_+,
 \qquad
 r_\lambda=e^\lambda\prod_{\beta\in\mathcal B'}
                      \beta^{E_\beta(\lambda)},
\]
where $u_+=\max(u,0)$ and $e^\lambda e^\gamma=e^{\lambda+\gamma}$.
These expressions belong to $\C(x_{ia})[\Z^\Omega]$.  For a polynomial
$p\in P^{\mathcal W_\lambda}$, with
$\mathcal W_\lambda=\operatorname{Stab}_{\mathcal W}(\lambda)$, denote
the distinct-translate orbit sum by
\[
 \mathsf S_\lambda(p)=
 \sum_{w\in\mathcal W/\mathcal W_\lambda}w(p r_\lambda).
\]
The choice of representative of the coweight orbit does not affect the
space of such sums.

\begin{proposition}\label{prop:laurent-model}
Under the abelianization embedding of
Proposition~\ref{prop:minuscule}, the no-flavor Coulomb algebra is
\begin{equation}\label{eq:laurent-model}
 A_0=\left(\bigoplus_{\lambda\in\Z^\Omega}P r_\lambda\right)^{\mathcal W}.
\end{equation}
The identification preserves the coefficient ring $R_0=P^{\mathcal W}$
and the topological grading.
\end{proposition}

\begin{proof}
First, the direct sum on the right is an algebra.  Indeed,
\begin{equation}\label{eq:laurent-product}
 r_\alpha r_\gamma
 =r_{\alpha+\gamma}
  \prod_{\beta\in\mathcal B'}
   \beta^{E_\beta(\alpha)+E_\beta(\gamma)-E_\beta(\alpha+\gamma)},
\end{equation}
and every exponent is nonnegative.  The multiset $\mathcal B'$ is
$\mathcal W$-stable, so taking invariants gives a subalgebra.

For a positive or negative subset coweight, the weights of the deleted
loop at a vertex cancel the tangent Euler denominator in
Equation~\eqref{eq:minuscule}, up to a nonzero constant sign.  This is the
Grassmannian tangent-weight calculation: for a positive subset $I$,
the loop contributes
$\prod_{a\in I,\ b\notin I}(x_{ib}-x_{ia})$, while the tangent weights
are these roots, up to their common sign convention.  For a negative
subset the roots are reversed.  This is the local Euler-factor calculation
associated with minuscule flag localization; compare the pushforward
identity in \cite[Lemma~2.10]{Weekes19}.
Proposition~\ref{prop:minuscule} therefore
places every generator of $A_0$ in the right-hand side of
Equation~\eqref{eq:laurent-model}.  Conversely, it places every
$\mathsf S_\alpha(p)$ for a positive or negative subset coweight
$\alpha$ in $A_0$, since a constant sign can be removed.

We prove that all the orbit sums $\mathsf S_\lambda(p)$ belong to $A_0$
by strong induction on
\[
 \|\lambda\|^2=\sum_{(i,a)\in\Omega}\lambda_{ia}^2.
\]
The coefficient and positive or negative subset-coweight cases have
just been established.
If $\lambda$ has a positive coordinate, set
\[
 f(t)=1_{t>0},\qquad g(t)=t-1_{t>0};
\]
otherwise set
\[
 f(t)=-1_{t<0},\qquad g(t)=t+1_{t<0}.
\]
In either case define $\alpha=f(\lambda)$ and $\gamma=g(\lambda)$
coordinatewise.  Outside the base cases both $\alpha$ and $\gamma$
have squared norm strictly smaller than that of $\lambda$.
Both functions are nondecreasing on $\Z$, and their values have
compatible signs relative to zero.  Thus, for every pair of slots
$d,e$ and every slot $d$,
\[
 (\alpha_d-\alpha_e)(\gamma_d-\gamma_e)\geq0,
 \qquad \alpha_d\gamma_d\geq0.
\]
These statements apply across different colors as well.  Consequently
every exponent in Equation~\eqref{eq:laurent-product} is zero for this
pair, and $r_\alpha r_\gamma=r_\lambda$.

Order the slots at each vertex by their $\lambda$-coordinates.  Then
$\alpha$ and $\gamma$ are ordered in the same way.  The rearrangement
inequality gives, for every relative matching $w\in\mathcal W$,
\[
 \|\alpha+w\gamma\|^2
 \leq\|\alpha+\gamma\|^2=\|\lambda\|^2.
\]
In the equality case, the resulting sum lies in the orbit of $\lambda$.
To see the last assertion directly, sort the matched entries within
each color.  Swapping an inverted pair strictly increases the scalar
product unless one of the two corresponding pairs is tied; swaps
involving ties leave the multiset of sums unchanged.  Hence every
maximal matching has the same multiset of sums as the ordered matching.
All other orbits occurring in a product
$\mathsf S_\alpha(p_\alpha)\mathsf S_\gamma(p_\gamma)$ therefore have
strictly smaller squared norm.

We now verify that the leading coefficient can be an arbitrary
$\mathcal W_\lambda$-invariant dressing.  At the fixed representative
$\lambda$, the decomposition into translates of $\alpha$ and $\gamma$
is unique.  For example, in the positive case, a different assignment
of the $1$-entries of $\alpha$ would assign $1$ to some coordinate
$\lambda_d\leq0$ and $0$ to some coordinate $\lambda_e>0$ of the same
color.  The corresponding $\gamma$-entries would satisfy
$\gamma'_d=\lambda_d-1<\lambda_e=\gamma'_e$, which is a strict
inversion relative to $\alpha'$.  This contradicts maximality.  The
negative case would similarly assign $0$ to a negative coordinate
$\lambda_d$ and $-1$ to a nonnegative coordinate $\lambda_e$ of the
same color; then $\gamma'_d=\lambda_d<\lambda_e+1=\gamma'_e$ is again
a strict inversion relative to $\alpha'$.  Thus the
coefficient of $r_\lambda$ in that product is precisely
$p_\alpha p_\gamma$.

Within each color, the level blocks of $\alpha$ and $\gamma$ are
intervals in the chosen ordering.  Their combined cuts are exactly the
cuts of the level blocks of $\lambda$: two distinct values cannot have
both the same $f$-value and the same $g$-value, since $f(t)+g(t)=t$.
For every $k\geq1$, the power sum on each $\lambda$-block is therefore
a linear combination of power sums on $\alpha$-blocks and
$\gamma$-blocks.  Explicitly, obtain prefix power sums at all cuts
from whichever block partition contains the cut, and subtract
successive prefixes.  Power sums generate each symmetric-polynomial
ring over $\C$, so
\[
 P^{\mathcal W_\lambda}
 =\bigl\langle P^{\mathcal W_\alpha},
                 P^{\mathcal W_\gamma}\bigr\rangle.
\]
As the rings commute, every element of the ring on the right is a
finite sum of products $p_\alpha p_\gamma$ from the two indicated
subrings.  Such products of dressed orbit sums realize any prescribed
leading dressing.  The other terms have polynomial dressings by
Equation~\eqref{eq:laurent-product}; their dressings are invariant under
the stabilizers of their coweights because the full product is
$\mathcal W$-invariant.  The induction hypothesis removes all those
lower terms.  This proves the reverse inclusion, without inverting any
polynomial, and completes the proof.
\end{proof}

\subsection{Multipartitions and cut ideals}

For the rest of this Section, $q$ records half the cohomological degree:
each polynomial variable has cohomological degree two and contributes $q$.
Thus
$\Hilb(M;q)=\sum_d\dim(M^{2d})q^d$ for the evenly graded vector spaces
under consideration.

Let $\mathcal M(\mathbf v)$ be the finite set of arrays
\[
 m=(m_{i,s})_{i\in Q_0,\ s\geq1},\qquad
 m_{i,s}\in\Nbb,\qquad \sum_{s\geq1}s m_{i,s}=v_i.
\]
For $m\in\mathcal M(\mathbf v)$, set
\[
 l_i=\sum_{s\geq1}m_{i,s},\qquad
 \mathbf l=(l_i)_i,\qquad |l|=\sum_i l_i,
\]
and
\[
 D_m=\frac12\sum_{i,j}\sum_{s,t\geq1}
                c_{ij}\min(s,t)m_{i,s}m_{j,t}.
\]
Let $\Omega_m^+$ consist of $m_{i,s}$ distinguished copies of each
type $(i,s)$.  These are the \emph{positive slots}.  For
$d\in\Omega_m^+$ denote its color by $i(d)$ and its size by $s(d)$.
Define
\[
 P_m=\C[x_d:d\in\Omega_m^+],\qquad
 \mathcal W_m=\prod_{i,s}S_{m_{i,s}},\qquad
 b_{de}=c_{i(d)i(e)}\min(s(d),s(e)).
\]
The group $\mathcal W_m$ permutes copies of the same type.  For every
nonempty subset $S\subseteq\Omega_m^+$, put
\begin{equation}\label{eq:cut-polynomial}
 q_S=
 \prod_{d\in S}x_d^{w_{i(d)}}
 \prod_{\substack{d\in S\\ e\in\Omega_m^+\setminus S}}
                    (x_d-x_e)^{b_{de}}.
\end{equation}
Let $J_m\subseteq P_m$ be the ideal generated by these cut polynomials,
and set
\[
 h_m(q)=\Hilb\bigl((P_m/J_m)^{\mathcal W_m};q\bigr).
\]
The ideal is homogeneous and $\mathcal W_m$-stable.  For the empty
array, these conventions give $P_m=\C$, $J_m=0$, and $h_m(q)=1$.

These cut products have a familiar graphical interpretation.
Adjoin a framing vertex $0$, put $w_{i(d)}$ edges between $0$ and each
positive slot $d$, and put $b_{de}$ edges between each unordered pair
of distinct slots $d,e$.
When this multigraph is connected, $J_m$ is the cocircuit ideal of its
reduced incidence configuration, with the row for $0$ deleted.
A minimal disconnecting cut contributes the product of its edge forms,
$x_d$ for framing edges and $x_d-x_e$ for slot edges; every cut product
is divisible by a minimal-cut product.
Such ideals occur in toric hyperk\"ahler cohomology
\cite[Section~7, Equation~(37)]{HauselSturmfels02} and in
permutation-invariant quiver constructions
\cite[Sections~2 and 4.4]{ReinekeRhoadesTewari23}.
A component disjoint from $0$ instead gives a unit cut and zero quotient.
Section~\ref{sec:counting} gives a direct reachability calculation that
retains the $\mathcal W_m$-action and yields the recurrence needed here.

\begin{proposition}\label{prop:filtration-bound}
The quotient $B_0$ for the full topological torus satisfies the
coefficientwise graded bound
\begin{equation}\label{eq:filtration-bound}
 \Hilb(B_0;q)\ \leq
 \sum_{m\in\mathcal M(\mathbf v)}
 q^{C(\mathbf v)-D_m+\mathbf w\cdot(\mathbf v-\mathbf l)}h_m(q).
\end{equation}
Every exponent displayed in the sum is a nonnegative integer.
\end{proposition}

\begin{proof}
Choose, for each color, an identification of its $l_i$ positive slots
with $l_i$ of its $v_i$ gauge slots.  On the remaining $v_i-l_i$ slots
put $s(d)=0$; call these the \emph{zero slots} and write their set as
$\Omega_m^0$.  Define a coweight by
\[
 (\lambda_m)_d=s(d)-1\qquad(d\in\Omega).
\]
The word ``positive'' refers to the part size, so at a rank-three vertex
the sizes $(2,1,0)$ give the coweight $(1,0,-1)$.
The coweight $\lambda_m$ has topological degree zero, since the sum of its coordinates at
vertex $i$ is $\sum_s s m_{i,s}-v_i=0$.  Its stabilizer is
\[
 \mathcal W_{\lambda_m}
 =\mathcal W_m\times\prod_i S_{v_i-l_i}.
\]
The last factors permute the zero slots.

By the degree convention in Proposition~\ref{prop:minuscule}, the
half cohomological degree of $r_{\lambda_m}$ is the polynomial degree
of its Laurent coefficient.  For a remaining arrow $i\to j$, that
degree is
\[
 \sum_{a=1}^{v_i}\sum_{b=1}^{v_j}
       \max(s(ia)-s(jb),0)
 =v_iv_j-\sum_{a=1}^{v_i}\sum_{b=1}^{v_j}
                         \min(s(ia),s(jb)),
\]
because $\sum_a s(ia)=v_i$.  The framing contribution at $i$ is
$w_i\sum_{d\in\Omega_m^+,\ i(d)=i}(s(d)-1)=w_i(v_i-l_i)$.
Summing gives
\begin{equation}\label{eq:multipartition-shift}
 \delta_m:=\deg_q r_{\lambda_m}
 =C(\mathbf v)-D_m+\mathbf w\cdot(\mathbf v-\mathbf l).
\end{equation}
This also proves integrality and nonnegativity.  The formula includes
remaining loops: their ordered pairs of gauge slots give precisely
the diagonal term with $c_{ii}=2n_{ii}$.

For $-1\leq L\leq\sum_i v_i$, define $\mathcal F_{-1}=0$ and, for
$L\geq0$, let $\mathcal F_L\subseteq B_0$ be the sum of the images of
\[
 \mathsf S_{\lambda_m}\bigl(P^{\mathcal W_{\lambda_m}}\bigr),
 \qquad m\in\mathcal M(\mathbf v),\quad |l|\leq L.
\]
Here and below an orbit sum denotes its image in the quotient when it
is viewed in $B_0$.  This is a finite filtration by graded vector
subspaces.  The array $m_{i,1}=v_i$ has $\lambda_m=0$, so the last
stage contains the image of all coefficients.  Those coefficients
generate $B_0$, by Proposition~\ref{prop:minuscule}; hence
$\mathcal F_{\sum_i v_i}=B_0$.  Exhaustion is asserted in the quotient,
and requires no assertion about spanning the entire topological-degree-zero
part of $A_0$ before quotienting.
For later use denote the symbol map of each orbit type by
\[
 \Phi_m:P^{\mathcal W_{\lambda_m}}
 \longrightarrow\mathcal F_{|l|}/\mathcal F_{|l|-1},\qquad
 p\longmapsto[\mathsf S_{\lambda_m}(p)].
\]
It adds the half-degree shift $\delta_m$.

We first eliminate dressings in the zero slots.  For each color $i$
and $k\geq1$, define an ambient Laurent-ring derivation by
\[
 D_{i,k}(p e^\lambda)
 =p e^\lambda\sum_a\lambda_{ia}x_{ia}^k,
 \qquad D_{i,k}(p)=0\quad(p\in\C(x_{ia})).
\]
Additivity of the displayed multiplier in $\lambda$ proves the
Leibniz rule.  This derivation commutes with $\mathcal W$, preserves
the algebra in Equation~\eqref{eq:laurent-model}, and preserves every
topological homogeneous summand.  It therefore preserves the ideal
defining $B_0$.  Its induced derivation on $B_0$ is zero, since it
kills the coefficient generators.  Applying it to an arbitrarily
dressed orbit sum gives
\begin{equation}\label{eq:zero-slot-power-sums}
 \mathsf S_{\lambda_m}\left(
 p\left[
 \sum_{\substack{d\in\Omega_m^+\\i(d)=i}}
                   (s(d)-1)x_d^k
 -\sum_{\substack{z\in\Omega_m^0\\i(z)=i}}x_z^k
 \right]\right)=0\quad\text{in }B_0
\end{equation}
for every $p\in P^{\mathcal W_{\lambda_m}}$.  The invariant dressing
ring is the tensor product of the symmetric-polynomial rings on the
positive type blocks and the zero blocks.  Power sums generate each
zero-block ring over $\C$.  Equation~\eqref{eq:zero-slot-power-sums},
including its arbitrary multiplier $p$, replaces each such power sum
by the indicated positive-slot polynomial.  Thus positive-slot
dressings alone span the image of this orbit type.  These replacements
preserve degree.

Fix now a nonempty $S\subseteq\Omega_m^+$, and split the array of
sizes on all gauge slots as $s=a+b$, by putting $b_d=s(d)$ on $S$ and
$b_d=0$ elsewhere.  Set
\[
 \alpha=a-1,\qquad \beta=b.
\]
The topological degree of $\beta$ is a nonzero vector with
nonnegative entries, so every product
\[
 \mathsf S_\alpha(p_a)\mathsf S_\beta(p_b)
\]
vanishes in $B_0$.  The dressings here may be arbitrary polynomials
invariant under their respective coweight stabilizers, by
Proposition~\ref{prop:laurent-model}.

To analyze the terms of this product, fix $a$ and match it with a
colorwise permutation $b'$ of $b$.  The resulting coweight is $t-1$,
where $t=a+b'$ has nonnegative entries and $\sum_{i(d)=i}t_d=v_i$.
It therefore determines another array in $\mathcal M(\mathbf v)$.
Writing $A=\operatorname{supp}(a)$ and
$B'=\operatorname{supp}(b')$, its length is
\[
 |\operatorname{supp}(t)|
 =|A|+|B'|-|A\cap B'|=|l|-|A\cap B'|.
\]
Every overlap term belongs to $\mathcal F_{|l|-1}$.  For a disjoint
matching, the positive sizes of $t$ are the multiset union of those
of $a$ and $b$ at each color.  Hence its array is exactly $m$; no
other orbit type occurs at length $|l|$.

For the chosen disjoint decomposition, the excess coefficient in
$r_\alpha r_\beta/r_{\lambda_m}$ is $q_S$ up to a nonzero constant
sign.  Here is the calculation for each weight.  For an arrow from
slot $d$ to slot $e$, its excess exponent is
\[
 (a_d-a_e)_+ +(b_d-b_e)_+
                  -(a_d+b_d-a_e-b_e)_+.
\]
This is zero unless one slot belongs to $\operatorname{supp}(a)$ and
the other to $S$.  In that case it is
$\min(s(d),s(e))$.  For the framing weight $-x_d$, the excess is
\[
 (a_d-1)_+ +(b_d)_+ -(a_d+b_d-1)_+,
\]
which equals one on $S$ and zero elsewhere.  These give exactly the
exponents in Equation~\eqref{eq:cut-polynomial}.  With the weight
conventions above, the sign is $(-1)^{\epsilon_S}$, where
\[
 \epsilon_S=
 \sum_{d\in S}w_{i(d)}+
 \sum_{\substack{d\in S\\e\in\Omega_m^+\setminus S}}
 n_{i(d)i(e)}\min(s(d),s(e)).
\]
This integer depends only on the numbers of selected and unselected
copies of each type.  In particular the sign is constant on the
averaging orbit considered next.

It remains to check the full span of the leading dressing relations.
Let
\[
 H=\mathcal W_{\lambda_m},\qquad
 H_S=\mathcal W_\alpha\cap\mathcal W_\beta.
\]
Then $H_S\subseteq H$.  Its blocks are the selected and unselected
copies of each positive type, together with the zero block at every
color.  All disjoint decompositions of the fixed representative
$\lambda_m$ into translates of $\alpha$ and $\beta$ form exactly the
set $H/H_S$.  The coefficient of its leading orbit sum in the product
above is $(-1)^{\epsilon_S}\operatorname{Tr}_{H_S}^H(q_Sp_ap_b)$.
Since the sign is nonzero and all other terms have lower length, we
obtain
\begin{equation}\label{eq:cut-leading-trace}
 \Phi_m\bigl(\operatorname{Tr}_{H_S}^H(q_Sp_ap_b)\bigr)=0,
 \qquad
 \operatorname{Tr}_{H_S}^H(f)=\sum_{h\in H/H_S}h(f).
\end{equation}

Choose $p_a$ to depend only on the positive-size blocks of $a$ and
$p_b$ only on the positive-size blocks of $b$.  Such choices are
allowed by their coweight stabilizers.  Their products span the
$H_S$-invariant polynomial ring on the positive slots, since that ring
is the tensor product of the symmetric-polynomial rings on the
selected and unselected copies of every type.  Moreover $q_S$ is
$H_S$-invariant.  For an arbitrary $f\in P_m$, Reynolds averaging first
over $H_S$ and then over $H$ consequently shows that
\[
 \operatorname{Av}_H(q_S f)
 =\operatorname{Av}_H\bigl(q_S\operatorname{Av}_{H_S}(f)\bigr)
\]
is a linear combination of the relations in
Equation~\eqref{eq:cut-leading-trace}.  All normalizing factors are
nonzero constants in $\C$.  Averaging the generators $q_S$ with
arbitrary multipliers spans $J_m\cap P_m^{\mathcal W_m}$, because
$J_m$ is stable under $\mathcal W_m$ and averaging is exact in
characteristic zero.  Therefore all invariant cut-ideal relations
annihilate the symbol of this orbit type.

Combining the zero-slot elimination with these cut relations gives,
for every $L$, a surjection of graded vector spaces
\begin{equation}\label{eq:filtration-symbol-surjection}
 \bigoplus_{\substack{m\in\mathcal M(\mathbf v)\\|l|=L}}
 (P_m/J_m)^{\mathcal W_m}[\delta_m]
 \twoheadrightarrow \mathcal F_L/\mathcal F_{L-1}.
\end{equation}
Here $[\delta_m]$ means that polynomial half degree $d$ is placed in
half cohomological degree $d+\delta_m$.  Exactness of invariants
identifies the source with the invariant dressing quotient just
obtained.  Summing its Hilbert-series bound over the finite exhaustive
filtration proves Equation~\eqref{eq:filtration-bound}.
\end{proof}

\begin{remark}\label{rem:filtration-edge-cases}
If there are no gauge slots, the unique empty array gives
$A_0=B_0=\C$ and equality in Equation~\eqref{eq:filtration-bound}.
If $\mathbf v\ne0$ and $\mathbf w=0$, the cut on all positive slots
is $q_{\Omega_m^+}=1$, so every summand on the right is zero.
Correspondingly, $e^{\mathbf 1}$ is an invertible monopole of nonzero
topological degree, and $B_0=0$.  The same observation applies to any
unframed connected component.  Full subsets, empty zero blocks, and
rank-one gauge factors require no change in the proof.
\end{remark}

The filtration has reduced the upper bound for $B_0$ to finitely many cut
quotients $(P_m/J_m)^{\mathcal W_m}$.  Section~\ref{sec:counting} evaluates
their generating series and compares the sum in
Equation~\eqref{eq:filtration-bound} with the cohomology at determinant
stability.  No algebra presentation of $B_0$ is inferred from the
vector-space surjections onto its filtration quotients.

\section{Cut ideals and the Betti-number bound}
\label{sec:counting}

We retain the notation of Section~\ref{sec:laurent}, including the assumption
that the original quiver has a loop at every vertex of positive dimension.
All Hilbert series in this section use half cohomological degree: each
variable $x_d$ has degree one.  We first compute the cut quotients and then
identify the right-hand side of Proposition~\ref{prop:filtration-bound} with
the ordinary Poincar\'e polynomial.
We write $v=\mathbf v$, $w=\mathbf w$, and $l=(l_i)_i$ for the dimension
vectors in this section.

For the counting identities, allow $m=(m_{i,s})_{i\in Q_0,\ s\geq1}$ to be any finite
array of nonnegative integers.  Its slots, cut ideal, and permutation group
are defined as before.  Set
\[
 \|m\|_i=\sum_{s\geq1}s m_{i,s},\qquad
 w(m)=\sum_{i,s}w_i m_{i,s},\qquad
 (q)_m=\prod_{i,s}\prod_{j=1}^{m_{i,s}}(1-q^j).
\]
For two such arrays, write
\[
 b(k,m')=\sum_{i,j,s,t}c_{ij}\min(s,t)k_{i,s}m'_{j,t}.
\]
Thus $D_m=\tfrac12 b(m,m)$.  Empty products equal one, and the empty array
has $P_0=\C$, $J_0=0$, and $h_0(q)=1$.

\begin{lemma}\label{lem:finite-cut-quotient}
For every finite array $m$, the algebra $P_m/J_m$ is finite dimensional.
Consequently $h_m(q)$ is a polynomial with nonnegative integer coefficients.
\end{lemma}

\begin{proof}
The assertion for the empty array is immediate.  Suppose that $x$ is a
common zero of the cut polynomials and that one of its coordinates has a
nonzero value $a$.  Take
\[
 S=\{d:x_d=a\}.
\]
Every framing factor in $q_S(x)$ is nonzero, and every difference across the
cut is nonzero because $S$ contains all slots with value $a$.  Thus
$q_S(x)\ne0$, a contradiction.  The common zero set of $J_m$ is therefore
contained in the origin.  By the Nullstellensatz, either $J_m=P_m$ or a
power of each variable belongs to $J_m$.  Both alternatives give a
finite-dimensional quotient.  Taking invariants under the finite group
$\mathcal W_m$ preserves finite dimensionality and the nonnegative grading.
\end{proof}

\begin{proposition}\label{prop:cut-recurrence}
For every finite array $m$,
\begin{equation}\label{eq:cut-recurrence}
 \frac{1}{(q)_m}
 =\sum_{k+m'=m}
   q^{w(k)+b(k,m')}\frac{h_{m'}(q)}{(q)_k}.
\end{equation}
\end{proposition}

\begin{proof}
We realize the cut products as Euler classes of incoming coordinates and
stratify by unreachable slots to obtain both the coefficient-ring
presentation and the recurrence.  Make a directed graph whose vertices
are the slots of $m$.
For each unordered pair of distinct slots $d,e$, put $b_{de}$ arrows from $d$ to $e$
and $b_{de}$ arrows from $e$ to $d$, with each ordered arrow counted once.
Add a fixed framing vertex and $w_{i(d)}$ arrows from it to $d$.
Assign a one-dimensional vector space to every slot and to the framing
vertex.  Let $E_m$ be the space of representations of this graph, and let
\[
 T_m=(\C^\times)^{\Omega_m^+}
\]
act by change of basis at the slots, leaving the framing line fixed.  The
class $x_d\in H_{T_m}^2(\mathrm{pt})$ is the first Chern class of the standard
character at slot $d$.

A slot is \emph{reachable} if there is a directed path from the framing
vertex to it on which every arrow coordinate is nonzero.  Let $U_m$ be the
open subset where all slots are reachable.  We will show that the natural
coefficient map gives
\begin{equation}\label{eq:cut-cohomology}
 H_{T_m}^*(U_m;\C)=P_m/J_m,
\end{equation}
and that this cohomology is concentrated in even degrees.  Reachability is
an open condition: it is a union of conditions requiring the coordinates
on finitely many paths to be nonzero.

We use induction on the number $n$ of slots.  For $n=0$, the representation
space and its reachable open are both a point, and both assertions hold.
For $0\leq r\leq n$, let $U_r\subset E_m$ be the open subset with at most
$r$ unreachable slots.  Thus $U_n=E_m$ and $U_0=U_m$.  For a subset $S$ of
$r\geq1$ slots, let $Z_S\subset U_r$ consist of representations whose
unreachable set is exactly $S$.

Let $L_S\subset E_m$ be the linear subspace on which all framing arrows
into $S$ and all arrows from the complement of $S$ into $S$ vanish.  These
equations force every slot of $S$ to be unreachable.  Conversely, within
$U_r$, unreachability of all $r$ slots in $S$ forces every complementary
slot to be reachable, so these incoming coordinates must vanish.  It follows that
\[
 Z_S=L_S\cap U_r.
\]
In particular, $Z_S$ is smooth and closed in $U_r$, and
$U_r\setminus U_{r-1}$ is the disjoint union of the $Z_S$ with $|S|=r$.

Let $k$ record the types of the slots in $S$, and put $m'=m-k$.  Restricting
to the complement identifies $Z_S$ with a vector bundle over $U_{m'}$.
Its remaining coordinates are the internal arrows on $S$ and the arrows
from $S$ to its complement, all of which are unrestricted.  Contracting
these vector-bundle fibers is equivariant.  The factor of $T_m$ acting at
the slots of $S$ acts trivially on the base, so the inductive hypothesis
gives
\begin{equation}\label{eq:cut-stratum-cohomology}
 H_{T_m}^*(Z_S)
 =H_{T_{m'}}^*(U_{m'})[x_d:d\in S]
 =(P_{m'}/J_{m'})[x_d:d\in S].
\end{equation}

The normal bundle is the representation formed by the incoming coordinates
set to zero.  Its complex rank and equivariant Euler class are
\begin{align}
 c_S&=w(k)+b(k,m'),\label{eq:cut-stratum-codimension}\\
 e_S&=\prod_{d\in S}x_d^{w_{i(d)}}
       \prod_{d\in S,\ e\notin S}(x_d-x_e)^{b_{de}}
       =q_S.\label{eq:cut-stratum-euler}
\end{align}
Only arrows entering $S$ vanish; this accounts for the exponent
$b(k,m')$, even though the graph has arrows in both directions.

Multiplication by $e_S$ on Equation~\eqref{eq:cut-stratum-cohomology} is
injective.  Indeed, view it as a polynomial in the variables indexed by
$S$, over the possibly nonreduced coefficient ring $P_{m'}/J_{m'}$.  Its
term of highest total degree in those variables is the single monomial
\[
 \prod_{d\in S}x_d^{w_{i(d)}+\sum_{e\notin S}b_{de}}
\]
with coefficient one.  Multiplication by this monomial is injective over
any coefficient ring, and all other terms have lower degree in the new
variables.  A highest-degree argument proves the claim, including the
case $e_S=1$.

Write $i_S:Z_S\hookrightarrow U_r$.  The self-intersection identity says
that $i_S^*i_{S,*}$ is multiplication by $e_S$.  Restriction to a different,
disjoint stratum kills the Gysin image.  Therefore the direct sum of the
Gysin maps is injective.  The equivariant Thom exact sequence, applied
successively from $U_n$ down to $U_0$, gives short exact sequences
\begin{equation}\label{eq:cut-thom-sequence}
 0\longrightarrow
 \bigoplus_{|S|=r}H_{T_m}^{*-2c_S}(Z_S)
 \xrightarrow{\ \sum i_{S,*}\ }H_{T_m}^*(U_r)
 \longrightarrow H_{T_m}^*(U_{r-1})\longrightarrow0.
\end{equation}
It also proves at each step that the cohomology is even and that
restriction from $P_m=H_{T_m}^*(E_m)$ is surjective.  These statements start
with $U_n=E_m$; evenness of the strata follows from the induction on $n$.

We also identify the kernel at each step.  By
Equation~\eqref{eq:cut-stratum-cohomology}, every class on $Z_S$ is the
restriction of a polynomial in the slot variables.  The incoming linear
equations cut out $Z_S$ regularly inside $U_r$, so its fundamental class
$i_{S,*}(1)$ is the restriction of the Euler product $q_S$.  The projection
formula now shows that the image of $i_{S,*}$ is the ideal generated by
that class.  Removing all strata therefore gives exactly the ideal
generated by the nonempty cuts, proving Equation~\eqref{eq:cut-cohomology}.

Finally, choose arrow labels compatibly with permutations of equal types.
Then $\mathcal W_m$ acts on the entire construction, permuting the slot
torus and its variables.  Complex orientations give equivariant Thom
classes without permutation signs.  Taking $\mathcal W_m$-invariants is
exact over $\C$.  Subsets $S$ with a fixed count vector $k$ form one orbit,
whose stabilizer is $\mathcal W_k\times\mathcal W_{m'}$.  Their total
invariant cohomology contribution in
Equation~\eqref{eq:cut-thom-sequence} is consequently
\[
 q^{c_S}
 \Hilb\bigl(\C[x_d:d\in S]^{\mathcal W_k};q\bigr)
 \Hilb\bigl((P_{m'}/J_{m'})^{\mathcal W_{m'}};q\bigr)
 =q^{w(k)+b(k,m')}\frac{h_{m'}(q)}{(q)_k}.
\]
There is one such contribution for each count vector $k$; taking invariants
of the sum over its subset orbit supplies no binomial multiplicity.
The all-reachable open supplies the term $k=0$.  Summing all strata and
using $\Hilb(P_m^{\mathcal W_m};q)=1/(q)_m$ proves
Equation~\eqref{eq:cut-recurrence}.
\end{proof}

\begin{corollary}\label{cor:cut-generating-ratio}
In the formal power series ring $\C(q)[[z_i:i\in Q_0]]$,
\begin{equation}\label{eq:cut-generating-ratio}
 \sum_m z^{\|m\|}q^{-D_m-w(m)}h_m(q)
 =
 \frac{\displaystyle\sum_m
       z^{\|m\|}q^{-D_m-w(m)}/(q)_m}
      {\displaystyle\sum_m z^{\|m\|}q^{-D_m}/(q)_m}.
\end{equation}
\end{corollary}

\begin{proof}
The quadratic form satisfies
\[
 D_{k+m'}=D_k+D_{m'}+b(k,m').
\]
This identity includes the diagonal entries of the type matrix: slots in
the two separate blocks are distinct even when their types agree.
Multiply Equation~\eqref{eq:cut-recurrence} by
$z^{\|m\|}q^{-D_m-w(m)}$ and sum over $m$.  Its summand factors as
\[
 \frac{z^{\|k\|}q^{-D_k}}{(q)_k}\,
 z^{\|m'\|}q^{-D_{m'}-w(m')}h_{m'}(q).
\]
For each monomial in the $z_i$, only finitely many arrays occur.  The
resulting convolution identity may therefore be divided by the first
factor's series, whose constant coefficient is one.  This is
Equation~\eqref{eq:cut-generating-ratio}.
\end{proof}

\subsection{Ordinary cohomology and Hausel's formula}

Let $\nu_{\det}=\prod_i\det_{V_i}$ denote determinant stability, and write
$X_{\det}(v,w)$ for the corresponding quiver variety.  This is a regular
choice: its pairing with every nonzero nonnegative dimension vector is
strictly positive.  Here is the corresponding GIT criterion, including
freeness.  Write $I_i:W_i\to V_i$ for the incoming framing maps.  In the
convention that semi-invariants satisfy $f(gx)=\nu_{\det}(g)^a f(x)$,
semistability is cyclicity: the smallest collection of subspaces invariant
under all doubled-quiver arrows and containing every $\im I_i$ is $V$.
If that collection is proper, a one-parameter subgroup with weight zero
on it and weight $-1$ on a chosen complement has a limit on the point and
negative determinant weight.  Every semi-invariant of positive determinant
weight must therefore vanish.  Conversely, at a cyclic point choose a
basis at each vertex from images of framing vectors under arrow paths.
The product of their basis determinants is a nonzero semi-invariant of
weight $\nu_{\det}$.  On the principal open where this product is nonzero,
normalizing these path bases gives a slice meeting each orbit once.  A
semistable orbit limit is cyclic, and one of its such principal opens also
meets the original orbit; the slice then shows that the limit belongs to
that orbit.  Thus these orbits are closed in the semistable locus.
Moreover, a stabilizer fixes all path images of framing vectors, and hence
is the identity.  With the opposite linearization convention the dual
argument gives cocyclicity.  The argument applies to the closed moment
fiber and allows loops throughout.  Put
\[
 P_{X_{\det}(v,w)}(q)
 =\sum_{j\geq0}\dim H^{2j}(X_{\det}(v,w);\C)q^j.
\]
Odd cohomology vanishes by Kirwan surjectivity, so this is the full ordinary
cohomology polynomial in half degrees.  The empty variety has polynomial
zero.

\begin{proposition}\label{prop:hausel-identification}
For a quiver with a loop at every vertex of positive dimension,
\begin{equation}\label{eq:hausel-cut-bound}
 P_{X_{\det}(v,w)}(q)
 =\sum_{\|m\|=v}
   q^{C(v)-D_m+w\cdot(v-l)}h_m(q).
\end{equation}
\end{proposition}

\begin{proof}
We first specify the moment-level comparison.  Give $N$ a Hermitian
metric invariant under $G_c=\prod_iU(V_i)$, and equip $T^*N$ with its flat
hyperk\"ahler structure.  The hyperk\"ahler moment map has one real and one
complex component.  Under the trace identification of the central real
dual with $\mathbb R^{Q_0}$, let $\eta=(1,\ldots,1)$ be the level corresponding to
the determinant character.  The GIT--moment-map correspondence identifies
$X_{\det}(v,w)$ with the compact-group quotient at real level $\eta$ and
complex level zero.  Rotation of the three moment-map components carries
this quotient to the one at real level zero and complex level $c\eta$,
for a common nonzero normalization constant $c$.  At real level zero,
affine Kempf--Ness identifies the latter quotient with
\[
 \mu_{\C}^{-1}(c\eta)/\!/G.
\]
These are identifications of underlying topological spaces; they do not
assert that rotation is holomorphic in the original complex structure.
The comparison uses the direct rotation on the flat representation
\cite[Section~3(D)]{HKLR87} and the topological affine Kempf--Ness
theorem recalled in Section~\ref{sec:preliminaries}.
The quiver GIT description is explained in \cite{Nakajima94}; the
flat-representation argument just given includes loops.
Multiplying the cotangent coordinate by $c^{-1}$
is an algebraic isomorphism from this affine quotient to
\[
 M_\eta(v,w)=\mu_{\C}^{-1}
       ((\id_{V_i})_i)/\!/G,
\]
because the cotangent moment map is linear in that coordinate.  Thus
$M_\eta(v,w)$ and $X_{\det}(v,w)$ have the same ordinary cohomology,
including simultaneous emptiness.

Hausel's formula for $M_\eta(v,w)$ explicitly allows loops and multiple
arrows \cite[Equation~(9)]{Hausel06}.  We record its grading and all factors
needed here.  For partitions $\lambda^i$ and $\lambda^j$, write
\[
 \langle\lambda^i,\lambda^j\rangle
 =\sum_{a,b}\min(\lambda_a^i,\lambda_b^j),\qquad
 E_Q(\lambda)
 =\sum_{a:i\to j\in Q_1}\langle\lambda^i,\lambda^j\rangle
   -\sum_i\langle\lambda^i,\lambda^i\rangle.
\]
Let $m(\lambda)_{i,s}$ be the number of parts of size $s$ in $\lambda^i$,
and set
\[
 \delta(v,w)=\sum_{a:i\to j\in Q_1}v_iv_j
                  -\sum_i v_i^2+w\cdot v.
\]
The nonempty smooth variety has complex dimension $2\delta(v,w)$.
With $q=t^2$, where $t$ is the ordinary cohomological variable in Hausel's
formula, that formula reads
\begin{equation}\label{eq:hausel-ordinary}
 \sum_v q^{-\delta(v,w)}P_{M_\eta(v,w)}(q)z^v
 =
 \frac{\displaystyle\sum_{\lambda}
       z^{|\lambda|}
       q^{-E_Q(\lambda)-\sum_i\langle\lambda^i,(1^{w_i})\rangle}
       /(q)_{m(\lambda)}}
      {\displaystyle\sum_{\lambda}
       z^{|\lambda|}q^{-E_Q(\lambda)}/(q)_{m(\lambda)}}.
\end{equation}
Here $\lambda$ ranges over all multipartitions, including the empty one.
In particular, the left side uses ordinary cohomology.  Its normalization
can also be checked by Poincar\'e duality: if $u$ is the arithmetic
variable, a smooth variety of complex dimension $2\delta$ satisfies
$P_c(u)=u^{2\delta}P(u^{-1})$.  The normalized count $u^{-\delta}P_c(u)$
therefore becomes $q^{-\delta}P(q)$ after setting $u=q^{-1}$, as in
Equation~\eqref{eq:hausel-ordinary}.

Now write $m=m(\lambda)$.  Removing one loop at each vertex cancels the
vertex self-pairing in $E_Q$, and therefore
\[
 E_Q(\lambda)=D_m,
 \qquad
 \sum_i\langle\lambda^i,(1^{w_i})\rangle
 =\sum_i w_i l_i=w(m),
 \qquad
 \delta(v,w)=C(v)+w\cdot v.
\]
Consequently the right side of Equation~\eqref{eq:hausel-ordinary} is
exactly the ratio in Equation~\eqref{eq:cut-generating-ratio}.
Taking its $z^v$ coefficient and multiplying by $q^{C(v)+w\cdot v}$ proves
Equation~\eqref{eq:hausel-cut-bound}.

The right side is a finite polynomial before this identification:
Lemma~\ref{lem:finite-cut-quotient} makes each $h_m$ finite, there are only
finitely many multipartitions of $v$, and the shifts are nonnegative
integers.  Explicitly,
\[
 C(v)-D_m
 =\tfrac12\sum_{i,j,s,t}c_{ij}
          \bigl(st-\min(s,t)\bigr)m_{i,s}m_{j,t}\geq0,
 \qquad
 w\cdot(v-l)=\sum_{i,s}w_i(s-1)m_{i,s}\geq0.
\]
The diagonal entries $c_{ii}$ are even and the off-diagonal terms pair,
which also verifies integrality.
\end{proof}

\begin{remark}\label{rem:counting-empty-cases}
When $v=0$, only the empty multipartition occurs and the polynomial is one.
When $w=0$ and $v\ne0$, the cut consisting of all positive slots has
$q_S=1$, so every summand is zero.  Equivalently, the numerator and
denominator of Equation~\eqref{eq:cut-generating-ratio} coincide.  This
agrees with emptiness at determinant stability: the diagonal gauge scalar
acts trivially on the fields and nontrivially on the linearization.  An
unframed connected component of the supported graph gives the same unit
cut.  Thus neither the recurrence nor the Betti comparison omits these
degenerate cases.
\end{remark}

\begin{corollary}\label{cor:loop-isomorphism}
Suppose that the quiver has a loop at every vertex of positive dimension.
At determinant stability, the coefficient map induces a canonical
isomorphism of graded $R_0$-algebras
\[
 B_0\xrightarrow{\ \sim\ }H^*(X_{\det}(v,w);\C).
\]
\end{corollary}

\begin{proof}
Proposition~\ref{prop:kernel-inclusion} gives the displayed graded
$R_0$-algebra surjection.  Proposition~\ref{prop:filtration-bound} and
Proposition~\ref{prop:hausel-identification} bound the dimension of each
graded piece of its source from above by that of its target.  Surjectivity
gives the reverse inequality.  Its kernel therefore vanishes in every
degree.  Since the bijection is the original coefficient-compatible
algebra map, this is an isomorphism over $R_0$ and an equality of the two
coefficient ideals, with the quotient rings retaining nilpotents.  The same argument
includes the zero ring in the empty case.
\end{proof}

\section{Flavor parameters, loop removal, and stability}
\label{sec:deformations}

Throughout this section, $B_F$ denotes the quotient by the full topological
torus.  In particular, it is independent of the stability character.
We first prove that regular stabilities have the same Betti numbers.
We then lift a no-flavor comparison to arbitrary flavor parameters and
use an added-loop parameter to remove the loop hypothesis at determinant
stability.  These three inputs combine in Theorem~\ref{thm:full-torus}.

\begin{lemma}\label{lem:betti-independence}
Fix a quiver and dimension and framing vectors.  The ordinary Betti numbers
of $X_\theta=\mu^{-1}(0)/\!/_{\theta}G$ are independent of the regular
character $\theta$, where regular means that every semistable point is stable
and the action on the semistable locus is free.  Empty quotients are included.
\end{lemma}

\begin{proof}
By Lemma~\ref{lem:local-stability}, each regular character has a real open
ball in which every rational parameter gives the same stable and semistable
loci, geometric quotient, and principal gauge bundle.  We may therefore
perturb within this ball without changing the quotient.

Here is a comparison of such neighborhoods for different regular
characters.  Write $G_c=\prod_i U(V_i)$ for a compact form and
$\mathfrak c=(\mathfrak g^*)^G$ for the space of complex central moment
parameters.  Equip $T^*N$ with its standard flat hyperk\"ahler structure.
The character space identifies with a real form of $\mathfrak c$.
For a rational parameter, clear its denominators and then dilate the
flat representation by the inverse square root of that common denominator;
the quadratic moment maps thereby return to level $\theta$.
Kempf--Ness identifies $X_\theta$ with the quotient by $G_c$ at real moment
level $\theta$ and complex moment level zero.  Quaternionic rotation commutes
with $G_c$ and rotates the three moment maps.  Thus it identifies this space
with the quotient at real level zero and complex level $c\theta$, where
$c\ne0$ is one fixed scalar determined by the normalization of the moment
maps.  Applying affine Kempf--Ness at this latter level gives a homeomorphism
\begin{equation}\label{eq:rotation-affine-fiber}
 X_\theta\ \simeq\ \mu^{-1}(c\theta)/\!/G .
\end{equation}
This uses the flat hyperk\"ahler quotient construction
\cite[Section~3(D)]{HKLR87} and the topological affine Kempf--Ness
identification recalled in Section~\ref{sec:preliminaries}; the
quiver-variety description appears in
\cite{Nakajima94}.  The rotation here acts directly on the flat
representation and is valid with loops.  It requires no properness of a
moment-map family.  One can normalize $c=1$; alternatively a linear dilation
of $T^*N$ identifies any two common nonzero scalar multiples of the complex
level.

Consider the affine family
\[
 \pi:\mathcal Z=
 \Spec\C[\{(x,\zeta)\in T^*N\times\mathfrak c:
                       \mu(x)=\zeta\}]^G
 \longrightarrow\mathfrak c.
\]
Its fiber at $\zeta$ is $\mu^{-1}(\zeta)/\!/G$.
To check this without a flatness assumption, let $A$ be the coordinate ring
inside the brackets.  The ideal specifying $\zeta$ is generated by invariant
parameter functions.  Exactness of invariants and the Reynolds operator
give
\[
 (A/\mathfrak m_\zeta A)^G
       =A^G/\mathfrak m_\zeta A^G.
\]
The morphism $\pi$ is of finite type.  Constructibility of compactly
supported direct image and its fiber-stalk description
\cite[Theorem~2.5(b)]{MaximSchurmann21}
\cite[Theorem~5.10]{SchnurerSoergel16}
therefore imply that $R\pi_!\C_{\mathcal Z}$ has constructible cohomology
sheaves with stalks $H_c^*(\pi^{-1}(\zeta);\C)$.  On a common
dense Zariski open subset $\Omega\subset\mathfrak c$, their ranks, hence all
compactly supported Betti numbers of the fibers, are constant.  We may take
$\Omega$ connected.  Every real open ball in the character space is Zariski
dense after multiplication by $c$, and its rational points are dense in the
usual topology.  Each regular character can thus be perturbed, without
changing its quotient, so that $c\theta\in\Omega$.

By Equation~\eqref{eq:rotation-affine-fiber}, all regular quotients have the same
compactly supported Betti numbers.  A nonempty such quotient is smooth of
complex dimension $2(\dim N-\dim G)$: freeness makes the moment map a
submersion along its stable locus.  Poincar\'e duality now gives equality of
ordinary Betti numbers.  If one regular quotient is empty, all compactly
supported Betti numbers vanish for every regular quotient; a nonempty smooth
quotient would have nonzero top compactly supported cohomology, a
contradiction.  This also covers nonfaithful gauge actions and unframed
components: the generic central fiber is allowed to be empty.  If $G$ is
trivial, the parameter space is a point and the assertion is immediate.
\end{proof}

\begin{lemma}\label{lem:flavor-lift}
For a fixed regular stability, suppose that the coefficient map induces
$B_0\simeq H^*(X)$.  Then, for every commuting flavor torus $F$, the map of
Proposition~\ref{prop:kernel-inclusion} is an isomorphism
\[
 B_F\xrightarrow{\ \sim\ }H_F^*(X)
\]
of graded $R$-algebras.
\end{lemma}

\begin{proof}
Set $S=\C[\mathfrak f]$ and let $S_+$ be its positive-degree ideal.
Kirwan surjectivity implies $H^{\mathrm{odd}}(X)=0$.
The Borel spectral sequence for $X\to EF\times_F X\to BF$ has second page
$H^*(BF)\otimes H^*(X)$, concentrated in even total degrees.  It degenerates,
and its edge map to $H^*(X)$ is surjective.  Lifting a finite homogeneous
basis of $H^*(X)$ and applying Leray--Hirsch shows that $H_F^*(X)$ is a free
$S$-module with
\[
 H_F^*(X)/S_+H_F^*(X)=H^*(X).
\]
This remains valid for $X=\varnothing$, with zero modules.

Let $J$ be the kernel of the surjection
$B_F\twoheadrightarrow H_F^*(X)$ from
Proposition~\ref{prop:kernel-inclusion}.  Freeness of the target makes
$\operatorname{Tor}_1^S(H_F^*(X),S/S_+)$ vanish.  Tensoring its defining exact
sequence with $S/S_+$ consequently gives
\[
 0\longrightarrow J/S_+J\longrightarrow B_F/S_+B_F
       \longrightarrow H^*(X)\longrightarrow0.
\]
Lemma~\ref{lem:flavor-base-change} identifies the middle term and map with
$B_0\xrightarrow{\sim}H^*(X)$; hence $J=S_+J$.
The module $J$ is graded and bounded below, since $B_F$ is a quotient of the
nonnegatively graded ring $R$.  If $J$ had a nonzero homogeneous element of
least degree, the identity $J=S_+J$ would express it using elements of smaller
degree, which is impossible.  Thus $J=0$.
This argument uses positivity only for the coefficient quotient $B_F$, and
does not assume that the Coulomb algebra itself is nonnegatively graded.
\end{proof}

\begin{proposition}\label{prop:loop-removal}
For every finite quiver and every dimension and framing vector, at the
determinant stability used in Corollary~\ref{cor:loop-isomorphism}, the
coefficient map induces an isomorphism of graded $R_0$-algebras
\[
 B_0\xrightarrow{\ \sim\ }H^*(X).
\]
\end{proposition}

\begin{proof}
Discard zero-rank gauge factors.  Add one loop at each remaining vertex and
write
\[
 N^+=N\oplus\bigoplus_i\End(V_i).
\]
Let $H=\C^\times$ act trivially on $N$ and with weight one on every added
loop.  Put $H^*_H(\mathrm{pt})=\C[z]$, with $\deg z=2$, and use a superscript
$+$ for the enlarged theory.  At determinant stability,
Corollary~\ref{cor:loop-isomorphism} and Lemma~\ref{lem:flavor-lift} give
\begin{equation}\label{eq:enlarged-flavor-isomorphism}
 B_H^+\xrightarrow{\ \sim\ }H_H^*(X^+)
 \qquad\text{as graded }R_0[z]\text{-algebras}.
\end{equation}
We compare this with the original theory endowed with the trivial $H$-action.
Let $\mathfrak m_0\subset R_0$ be the maximal ideal of the graded origin and
set
\begin{equation}\label{eq:loop-local-ring}
 \Lambda=R_0[z]_{(\mathfrak m_0,z-1)}.
\end{equation}
The comparison over $\Lambda$ is ungraded in cohomological degree but
retains topological charge.  The equality of homogeneous ideals in $R_0$
obtained below will restore the cohomological grading.

\emph{The Coulomb comparison.}
Use the common rational Laurent realization of
Proposition~\ref{prop:minuscule}, and put $P=\C[x_{ia}]$.
For a positive or negative minuscule coweight $\lambda$, the added numerator
in Equation~\eqref{eq:minuscule} is
\[
 E_\lambda(x,z)=
   \prod_i\prod_{a,b=1}^{v_i}
     (z+x_{ib}-x_{ia})^{(\lambda_{ia}-\lambda_{ib})_+}.
\]
It is invariant under the stabilizer $\mathcal W_\lambda$ and has value one
at $(x,z)=(0,1)$.  Its inverse belongs to
$P[z]^{\mathcal W_\lambda}\otimes_{R_0[z]}\Lambda$: indeed,
\[
 \Delta_\lambda=\prod_{w\in\mathcal W}wE_\lambda
       \in R_0[z],\qquad \Delta_\lambda(0,1)=1,
\]
and $E_\lambda^{-1}=(\Delta_\lambda/E_\lambda)/\Delta_\lambda$.
If $M(\lambda,p)$ denotes an original dressed generator, the enlarged
generator with dressing $p$ is $M(\lambda,pE_\lambda)$.  Conversely,
dressing by $p/E_\lambda$ recovers $M(\lambda,p)$ after localization.
Generation therefore identifies the two localized Coulomb algebras, with
their coefficient maps and their full topological gradings.  The elements
inverted in Equation~\eqref{eq:loop-local-ring} have topological degree zero.  Thus
localization commutes with quotienting by all nonzero topological degrees,
and gives
\begin{equation}\label{eq:loop-coulomb-localization}
 B_H^+\otimes_{R_0[z]}\Lambda
    \simeq (B_0\otimes_\C\C[z])\otimes_{R_0[z]}\Lambda.
\end{equation}

\emph{The cohomological comparison.}
Let $D^+\subset T^*N^+$ be the stable zero moment locus.  An $H$-fixed
point of the quotient may be fixed only after a compensating gauge
transformation.  We therefore localize the actual $H$-fixed locus of this
prequotient, retaining $G\times H$ equivariance, and first restrict to
$T\times H$.  Brion's subgroup localization theorem
\cite[Lemma~4]{Brion98} says that, for a torus space embedded equivariantly in
a finite-dimensional representation, restriction to the fixed locus of a
subtorus becomes an isomorphism after inverting finitely many characters
nontrivial on that subtorus.  Apply it to the compact forms of $T\times H$
and $H$ and to the invariant locally closed subset $D^+\subset T^*N^+$.
Every character nontrivial on $H$ is nonzero at the central Lie algebra
element $(0,1)$.  The restriction
\[
 H^*_{T\times H}(D^+)\longrightarrow H^*_{T\times H}((D^+)^H)
\]
is consequently an isomorphism there.  Taking $\mathcal W$-invariants
recovers $G\times H$ equivariance.  This is compatible with localization
over $\Lambda$: the point is $\mathcal W$-fixed, and the product of the Weyl
translates of any denominator is an invariant denominator still nonzero
there.  This is also the central-parameter case of the prequotient
localization in \cite[Proposition~2.3]{DK25}.

Both the added loops and their cotangent partners have nonzero $H$-weights,
so $(D^+)^H$ sets all these coordinates to zero.  The enlarged moment map
then restricts to the original moment map.  Stability restricts to the
original stability as well: under a $G$-equivariant closed inclusion, the
Hilbert--Mumford test of a point and of every one-parameter subgroup is
unchanged.  Equivalently, semi-invariants lift by reductivity, and orbit
closures remain in the closed subvariety.  Hence
\[
 (D^+)^H=D,
\]
with $H$ acting trivially on the right.  Freeness of the stable actions now
gives a coefficient-compatible isomorphism
\begin{equation}\label{eq:loop-higgs-localization}
 H_H^*(X^+)\otimes_{R_0[z]}\Lambda
    \simeq (H^*(X)\otimes_\C\C[z])\otimes_{R_0[z]}\Lambda.
\end{equation}
The localization theorem used here does not require compactness of $D^+$
or $X^+$.

\emph{Descent.}
Let $I_B$ and $I_X$ be the kernels of the coefficient maps from $R_0$ to
$B_0$ and $H^*(X)$, respectively.  Equations
\eqref{eq:enlarged-flavor-isomorphism},
\eqref{eq:loop-coulomb-localization}, and
\eqref{eq:loop-higgs-localization} identify the localized quotients over the
same ring $\Lambda$.  Therefore
$I_B\Lambda=I_X\Lambda$.  Setting $z=1$ uses the identity
\[
 \Lambda/(z-1)\Lambda=(R_0)_{\mathfrak m_0}
\]
and yields
$(I_B)_{\mathfrak m_0}=(I_X)_{\mathfrak m_0}$.
Both ideals are homogeneous.  If $h\in I_B$ is homogeneous, there is
$s\in R_0$ with $s(0)\ne0$ and $sh\in I_X$.  The component of degree
$\deg h$ of this relation is $s(0)h\in I_X$, whence $h\in I_X$.
Interchanging the ideals proves $I_B=I_X$ globally.  This is equality of
ideals and gives the asserted graded isomorphism of quotient rings, with
all nilpotents retained.  The proof applies unchanged to empty stable loci; if there are
no nonzero gauge ranks, the assertion is immediate.
\end{proof}

\begin{theorem}\label{thm:full-torus}
For every datum in the statement of Theorem~\ref{thm:main}, the coefficient
map induces a canonical isomorphism of graded $R$-algebras
\[
 B_F\xrightarrow{\ \sim\ }H_F^*(X).
\]
\end{theorem}

\begin{proof}
At determinant stability, Proposition~\ref{prop:loop-removal} identifies
$B_0$ with ordinary cohomology, so $B_0$ is finite-dimensional.  For any
other regular stability, Lemma~\ref{lem:betti-independence} gives the same
total Betti number.  The graded coefficient-compatible surjection
$B_0\twoheadrightarrow H^*(X)$ of
Proposition~\ref{prop:kernel-inclusion} is therefore an isomorphism.
Lemma~\ref{lem:flavor-lift} proves the assertion for every commuting flavor
torus, with the generality supplied by Lemma~\ref{lem:trivial-summands}.
If a quotient is empty, the same argument gives $B_0=B_F=0$.
Uniqueness follows from surjectivity of the coefficient map $R\to B_F$.
\end{proof}

\section{The cocharacter fixed scheme}\label{sec:cocharacters}

We now compare the fixed scheme for the chosen cocharacter with the fixed
scheme for the full topological torus.  This comparison uses only minuscule
generation, a root-support consequence of the loop-nilpotent Hall algebra,
and the regularity of the stability.  In particular, it does not use
Theorem~\ref{thm:full-torus} or assume coefficient surjectivity onto the
cocharacter fixed quotient.

If \(\nu\cdot\mathbf b\ne0\) for every
\(0<\mathbf b\leq\mathbf v\), Proposition~\ref{prop:minuscule} immediately
gives the comparison: every nontrivial positive or negative minuscule
generator has nonzero \(\nu\)-weight.  The argument below also covers
regular characters that lie on some of these dimension hyperplanes.

For example, take two disconnected vertices, no arrows, and
\(\mathbf v=\mathbf w=(1,1)\), with \(\nu=(1,-1)\).
Each Higgs factor is a point with free stable gauge action, so this
character is regular, although \(\nu\cdot(1,1)=0\).
The joint positive monopole is a product of monopoles of weights
\(1\) and \(-1\), and hence vanishes in the cocharacter fixed quotient:
weight zero alone does not imply survival.

To treat these characters, we adjoin a rank-one framing vertex.
If a nontrivial minuscule generator of the original Coulomb algebra
survived in its cocharacter fixed quotient, the corresponding class and
its complementary class would both survive in the enlarged cocharacter
fixed quotient.  Polynomial Hall-algebra generation would then decompose
both nonzero complementary dimension vectors into indecomposable dimension
vectors of weight zero for the extended character.
These two decompositions supply at least two blocks of a semistable direct
sum; after removing the one ineffective diagonal scalar, its stabilizer
remains positive-dimensional, contradicting regularity.

\subsection{Generation in indecomposable charges}

For a finite quiver \(\Gamma\), possibly with loops, write
\[
 \Delta_\Gamma
 =\{\,0\ne\boldsymbol\gamma\in\Nbb^{\Gamma_0}:
                  A_{\Gamma,\boldsymbol\gamma}(q)\ne0\,\}.
\]
Here \(A_{\Gamma,\boldsymbol\gamma}\) is the polynomial counting
absolutely indecomposable representations over finite fields introduced
in Kac's work \cite{Kac83}; polynomiality for quivers with self-loops
is explicitly recalled in \cite[Theorem~1(a) and Remark~4]{Kac19}.
We use \(\Delta_\Gamma\) only as the nonzero support of these polynomials,
without a loopless root-system interpretation.

The generation mechanism belongs to the cohomological Hall theory
initiated by Kontsevich--Soibelman \cite{KontsevichSoibelman11}.
Its BPS and PBW structure was developed by Davison--Meinhardt and
Davison \cite{DavisonMeinhardt20,Davison25AffineBPS}; the graded Kac
identity builds on Mozgovoy's generating-series formula and Davison's
cohomology of preprojective stacks
\cite{Mozgovoy11McKay,Davison23Preprojective}.
The difference-operator maps were developed in the work of
Finkelberg--Tsymbaliuk, Frassek--Tsymbaliuk, and Tsymbaliuk
\cite{FinkelbergTsymbaliuk19Integral,FrassekTsymbaliuk22,Tsymbaliuk23Difference}.
For arbitrary unframed quivers, Muthiah--Weekes relate Kac-polynomial
Hilbert series and shuffle algebras through limit quantized zastava
spaces \cite{MuthiahWeekes2026Zastavas}.
The zastava comparison provides related context.  Our inputs are the
polynomial-parameter loop-nilpotent package of Jindal--Negu\c{t}
\cite{JindalNegut26} and Tsymbaliuk's arbitrary-quiver operator maps
\cite[Theorem~5.7]{Tsymbaliuk23Difference}.
We first fix their product conventions and prove that the maps restrict
to polynomial coefficient rings.  This is needed before specializing
the independent parameters to the given flavor action.

Put
\[
 S=\Qbb[\hbar,u_a\mid a\in\Gamma_1],
 \qquad K=\operatorname{Frac}(S),
\]
with independent parameters.  Let \(\mathcal S_S^+\) be the polynomial
shuffle half of \cite[Definition~2.3]{JindalNegut26}, and denote its
product by \(*\).  It is graded by dimension vectors, and contains
\[
 e_{\mathbf n,g}
   =g(z_{ia})\prod_i\prod_{a,b=1}^{n_i}
                   (z_{ib}-z_{ia}+\hbar),
 \qquad
 g\in S[z_{ia}]^{\prod_iS_{n_i}}.
\]
For homogeneous elements of dimensions \(\mathbf n,\mathbf m\), use
the product
\begin{equation}\label{eq:hall-sign-product}
 E\star E'
   =(-1)^{\sum_{a:i\to j}n_jm_i}E*E'.
\end{equation}
Bilinearity of the exponent makes this product associative.
Let \(\mathcal S_S^-\) be its ordinary opposite algebra, with
\(f_{\mathbf n,g}\) denoting the same polynomial in that half.
A word changes under this twist by a sign depending only on the
dimensions of its factors; all homogeneous \(S\)-word spans are unchanged.

\begin{lemma}[Polynomial monopole maps]\label{lem:polynomial-monopole-maps}
For every unframed quiver \(\Gamma\) and dimension vector \(\mathbf d\),
the two halves above admit \(S\)-linear algebra maps to an algebra of
difference operators defined without inverting rotation or arrow parameters.
They send \(e_{\mathbf n,g}\) and \(f_{\mathbf n,g}\), for
\(\mathbf n\leq\mathbf d\), to nonzero constant multiples of the
positive and negative fundamental operators, respectively; they send
these elements to zero when \(\mathbf n\not\leq\mathbf d\).
After setting \(\hbar=0\), \(w_{ia}=x_{ia}\), and
\(u_a=-f_a\), where \(f_a\) is the prescribed flavor weight of arrow
\(a\), their images lie in \(A_{\Gamma,\mathbf d,F}\).
Dimension-vector degree becomes positive, respectively negative,
topological charge.  These maps preserve finite polynomial-coefficient
identities under this specialization.
\end{lemma}

\begin{proof}
We specify the operator algebra and the sign convention.
Let \(D_{ia}w_{jb}=(w_{jb}+\hbar\delta_{ij}\delta_{ab})D_{ia}\),
with all the \(D_{ia}\) commuting.  Allow coefficients in
\[
 S[w_{ia},(w_{ia}-w_{ib}+r\hbar)^{-1}:
                      a\ne b,\ r\in\Z].
\]
This ring is preserved by every shift \(D_{ia}^{\pm1}\).
For subsets \(A_i\subseteq\{1,\ldots,d_i\}\), write
\(w_A=(w_{ia})_{a\in A_i,i\in\Gamma_0}\).  Define
\begin{align*}
 E_{\mathbf n,g}
 &=\sum_{|A_i|=n_i}
   g(w_A)
   \frac{\displaystyle\prod_{a:i\to j}
             \prod_{\alpha\in A_i,\ \beta\notin A_j}
                (w_{i\alpha}-w_{j\beta}+u_a)}
        {\displaystyle\prod_i
             \prod_{\alpha\in A_i,\ \beta\notin A_i}
                (w_{i\alpha}-w_{i\beta})}
       \prod_{i,\alpha\in A_i}D_{i\alpha},\\
 F_{\mathbf n,g}
 &=\sum_{|A_i|=n_i}
   g(w_A-\hbar)
   \frac{\displaystyle\prod_{a:j\to i}
             \prod_{\alpha\in A_i,\ \beta\notin A_j}
                (w_{i\alpha}-w_{j\beta}-u_a)}
        {\displaystyle\prod_i
             \prod_{\alpha\in A_i,\ \beta\notin A_i}
                (w_{i\alpha}-w_{i\beta})}
       \prod_{i,\alpha\in A_i}D_{i\alpha}^{-1}.
\end{align*}
These are the fundamental formulas
\cite[Equations~(78)--(79)]{JindalNegut26}.
Let \(\mathcal C_S\) be the \(S\)-subalgebra generated by these operators
and the polynomials symmetric in \(w\) separately at each vertex.

We construct the maps first over \(K\), then prove that their polynomial
halves land in \(\mathcal C_S\).
Use Tsymbaliuk's maps on the full trigonometric shuffle algebra and its
opposite \cite[Theorem~5.7, Equations~(93) and~(96)]{Tsymbaliuk23Difference},
with no framing, and make the formal substitutions
\[
 q=e^{\varepsilon\hbar},\qquad
 t_a=e^{\varepsilon u_a},\qquad
 W_{ia}=e^{\varepsilon w_{ia}}.
\]
Here \(\varepsilon\) is an auxiliary formal variable, not \(\hbar\).
The multiplicative shift in the source satisfies
\(D^{\mathrm{trig}}_{ia}W_{ia}=q^{-1}W_{ia}D^{\mathrm{trig}}_{ia}\).
Thus its inverse becomes our positive shift \(D_{ia}\), while it
becomes \(D_{ia}^{-1}\) in the negative map.
The leading coefficient of the trigonometric shuffle kernel is
\[
 \left(\frac{x-\hbar}{x}\right)^{\delta_{ij}}
 \prod_{a:i\to j}[-(x+u_a)]
 \prod_{a:j\to i}(x+\hbar-u_a).
\]
It is precisely the kernel for \(\star\) in
Equation~\eqref{eq:hall-sign-product}: the incoming-arrow signs have
been removed from the kernel of \cite[Equation~(31)]{JindalNegut26}.
This distinction is necessary.  For the quiver consisting of one arrow
\(i\to j\), the untwisted product gives
\(e_{(1,0),1}*e_{(0,1),1}
   =-e_{(1,1),z_i-z_j+u_a}\), whereas the corresponding fundamental
operator product has a plus sign.

Here are the order calculations that justify taking leading coefficients,
including for loops and repeated gauge slots.
Put
\[
 \kappa(\mathbf n)=\sum_{a:i\to j}n_in_j,\qquad
 \ell_i=\#\{a:i\to i\},\qquad
 L_i^+=\sum_{a:i\to j}d_j-d_i+1,\qquad
 L_i^-=\sum_{a:j\to i}d_j-d_i+1.
\]
Lift a polynomial \(E\) of dimension \(\mathbf n\) to
\(\varepsilon^{\kappa(\mathbf n)}
 E((X_{ia}-1)/\varepsilon)\).
The cross-kernel order is
\(\kappa(\mathbf n+\mathbf m)-\kappa(\mathbf n)-\kappa(\mathbf m)\);
thus the leading source product is \(\star\).
In an operator summand, let \(m_{ia}\) be the length of the string
assigned to gauge slot \((i,a)\), so \(\sum_a m_{ia}=n_i\), and set
\(P_2=\sum_{i,a}\ell_i m_{ia}^2\).
The three prefactors in Tsymbaliuk's formulas have orders
\[
 \begin{array}{c|c}
 \text{factor}&\varepsilon\text{-order}\\ \hline
 \text{single-entry factors }Y\text{ or }Y'
     &L^\pm\cdot\mathbf n-\sum_i\ell_i n_i\\
 \text{inverse kernels within one string}
     &-P_2+\sum_i\ell_i n_i\\
 \text{factors between distinct slots}
     &-\kappa(\mathbf n)+P_2 .
 \end{array}
\]
Together with the lifted input, their total order is
\(L^\pm\cdot\mathbf n\), independently of the composition.
Multiply the operator image of the lift by
\(\varepsilon^{-L^\pm\cdot\mathbf n}\) and take its
\(\varepsilon^0\) coefficient to define the normalized leading map.
The remaining denominators have nonzero leading coefficients over
\(K(w)\).  There are only finitely many compositions.  Consequently
an input with zero normalized leading coefficient also has zero
normalized leading output.  The construction is independent of the lift.
Since \(L^\pm\cdot\mathbf n\) is additive, the normalized leading
maps preserve multiplication.  Multiplying them in dimension
\(\mathbf n\) by \(\hbar^{-|\mathbf n|}\) preserves that property.

The use of formal parameters in this argument requires no analytic
specialization theorem.  For real
\(\hbar<u_a<0\) and small positive \(\varepsilon\), the substitutions
satisfy Tsymbaliuk's parameter inequalities.
The monomial factors involving logarithms of torus parameters are
expanded near \(W=1\); they have leading coefficient one and rational
Taylor coefficients in \(\hbar,u_a,w_{ia}\).
Each coefficient of the multiplication identities is therefore a
rational identity valid on this nonempty real open set, hence valid
over \(K(w)\).  This justifies the formal leading maps, denoted
\(\Psi_+\) and \(\Psi_-\).

For the distinguished elements the maps are especially simple:
\begin{equation}\label{eq:corrected-monopole-images}
 \Psi_+(e_{\mathbf n,g})
   =(-1)^{L^+\cdot\mathbf n+\kappa(\mathbf n)}E_{\mathbf n,g},
 \qquad
 \Psi_-(f_{\mathbf n,g})=F_{\mathbf n,g}.
\end{equation}
Indeed, the same-color factor in \(e_{\mathbf n,g}\) vanishes on every
additive string of length at least two.  Only selected subsets of
distinct slots remain.  Its diagonal factors contribute
\(\hbar^{|\mathbf n|}\), which cancels the normalization
\(\hbar^{-|\mathbf n|}\); its off-diagonal factors cancel the
same-color denominators in the operator formula.  The factors for
arrows between two selected slots cancel as well.
In the positive formula the single-entry factors contribute
\((-1)^{\sum_i n_i(L_i^+-\ell_i)}\), and the internal-arrow
cancellations contribute
\((-1)^{\kappa(\mathbf n)-\sum_i\ell_i n_i}\), giving the stated sign.
In the opposite formula the selected variables are \(w_A-\hbar\);
the incoming-arrow factors cancel without these signs.
This proves both formulas, and proves that the image is zero when
\(\mathbf n\not\leq\mathbf d\).

By \cite[Theorem~2.4]{JindalNegut26}, the distinguished elements generate
\(\mathcal S_S^+\) over \(S\).  The twist preserves their word spans,
and the opposite statement follows by reversing words.
Equation~\eqref{eq:corrected-monopole-images} therefore proves that
both polynomial halves land in \(\mathcal C_S\).
Although we constructed the maps over \(K\), their restrictions are
thus defined over \(S\); no inverse of a parameter is needed in the
identities in \(\mathcal C_S\).

Finally specialize \(\hbar=0\), \(w_{ia}=x_{ia}\), and \(u_a=-f_a\).
Every allowed denominator becomes \(x_{ia}-x_{ib}\), with \(a\ne b\),
so this defines a homomorphism to
\(\operatorname{Frac}(P_F)[X_*(T)]\), sending each \(D_{ia}\)
to the monomial for its standard slot coweight.
For a selected-variable dressing \(g\), also denoting its specialization
by \(g\), direct comparison with Proposition~\ref{prop:minuscule} gives
\[
 M^+_{\mathbf n,g}
   =(-1)^{\sum_{a:i\to j}n_i(d_j-n_j)}
          E_{\mathbf n,g}\big|_{\mathrm{specialization}},
 \qquad
 M^-_{\mathbf n,g}
   =(-1)^{\sum_i n_i(d_i-n_i)}
          F_{\mathbf n,g}\big|_{\mathrm{specialization}}.
\]
The first sign comes from the outgoing matter factors; the second
comes from reversing the minuscule tangent weights.
The coefficient roots themselves are unchanged.
Thus every generator of \(\mathcal C_S\) specializes to an element
of \(A_{\Gamma,\mathbf d,F}\), and so does every image of the
polynomial halves.  Only this homomorphism is needed; no injectivity
or flatness after specialization is asserted.
\end{proof}

\begin{lemma}[Root-charge generation]\label{lem:root-charge-generation}
Let \(\Gamma\) be an unframed quiver with dimension vector \(\mathbf d\),
and let a flavor torus act by a character on each arrow.
Write \(A_{\Gamma,\mathbf d,F}\) for its commutative Coulomb algebra and
\(R_{\Gamma,\mathbf d,F}\) for its coefficient algebra.
For \(\epsilon\in\{+1,-1\}\), let
\(M_{\mathbf b,p}^{\epsilon}\) be a dressed minuscule class whose
coweight has \(b_i\) entries equal to \(\epsilon\) at vertex \(i\) and
all other entries zero, where \(0<\mathbf b\leq\mathbf d\).
There is a finite expression
\begin{equation}\label{eq:root-charge-words}
 M_{\mathbf b,p}^{\epsilon}
   =\sum_r c_r a_{r1}\cdots a_{rk_r},
 \qquad c_r\in R_{\Gamma,\mathbf d,F},
\end{equation}
in which every factor is topologically homogeneous and
\[
 \deg_{\mathrm{top}}a_{rh}=\epsilon\boldsymbol\gamma_{rh},
 \qquad
 \boldsymbol\gamma_{rh}\in\Delta_\Gamma,
 \qquad
 \sum_{h=1}^{k_r}\boldsymbol\gamma_{rh}=\mathbf b.
\]
The factors need not themselves be minuscule classes.
\end{lemma}

\begin{proof}
\emph{Polynomial Hall-algebra input.}
Use the independent polynomial ring \(S\) defined above.  This is
the equivariant coefficient ring $H^*(BT_{\mathrm{univ}};\Qbb)$ for
$T_{\mathrm{univ}}=(\C^\times)^{1+|\Gamma_1|}$, with independent
rotation and arrow characters.  It lies in the geometric parameter
setting of \cite[Equation~(30) and Remark~2.6]{JindalNegut26}; no
extension to arbitrary field-valued exponential parameters is needed.
On the
tripled quiver use weights \(-u_a\), \(u_a-\hbar\), and \(\hbar\)
on an original arrow, its reverse, and an added loop, respectively.
These choices satisfy all the genericity hypotheses of
\cite[\S2.2]{JindalNegut26}; in particular there is no nonzero
nonnegative integral relation among \(u_a,\hbar-u_a\), and
\(u_a\notin\Z\hbar\) for every original loop.

We use the following polynomial-form results of that source:
the loop-nilpotent Hall algebra identifies with its shuffle half
\cite[Proposition~3.7]{JindalNegut26};
its BPS descendants generate over \(S\)
\cite[Theorem~3.10, Equation~(109), and
Corollary~3.15]{JindalNegut26};
and the BPS comparison and Kac formula show that the nonzero
dimension-vector support of these generators is \(\Delta_\Gamma\)
\cite[Proposition~3.14 and \S3.10]{JindalNegut26}.
Here the Kac identity retains the full cohomological grading: the
coefficient of $t^{d/2}$ is $(-1)^d\dim H^d$.
A zero Kac polynomial therefore forces every graded piece of the
corresponding BPS space to vanish.  Proposition~3.14 of that source
transfers this vanishing to the loop-nilpotent BPS space.
The concrete twist in Equation~\eqref{eq:hall-sign-product} preserves
these homogeneous word spans.
Here polynomial form refers to equivariant parameters over \(\Qbb\),
not integral cohomology.

Apply the generation statement to \(e_{\mathbf b,g}\) or
\(f_{\mathbf b,g}\), and retain dimension-vector degree \(\mathbf b\).
All nonconstant Hall degrees are nonnegative, so the result is a
finite \(S\)-linear combination of products of BPS descendants with
degrees in \(\Delta_\Gamma\) summing to \(\mathbf b\).
Applying Lemma~\ref{lem:polynomial-monopole-maps} gives the corresponding
finite identity in \(\mathcal C_S\).
Its coefficients lie in \(S\); no parameter from \(S\) is inverted
in this identity.
Every factor has degree at most \(\mathbf b\leq\mathbf d\); no assertion
about vanishing of entire Hall components outside this bound is needed.

\emph{Specialization.}
Now extend scalars to \(\C\) and apply the specialization homomorphism
of Lemma~\ref{lem:polynomial-monopole-maps}.
Each finite identity becomes an identity in \(A_{\Gamma,\mathbf d,F}\).
The signs in the images of the distinguished elements are nonzero
constants and can be absorbed in its coefficients.
This operation is performed on the abstract polynomial-form algebra and
its map, not on the ambient polynomial representation of the shuffle
algebra.  An element whose ambient representative is divisible by
\(\hbar\) need not belong to \(\hbar\) times the polynomial shuffle
half, and is therefore not discarded before applying the map.

\emph{General dressings.}
It remains to allow an arbitrary stabilizer-invariant dressing \(p\).
At each vertex, the polynomial algebra symmetric separately in the
selected and complementary variables is generated by the polynomials
symmetric in all variables together with those symmetric in the
selected variables.  Indeed, each complementary power sum is the total
power sum minus the selected power sum.  Thus \(p\) is a finite sum
of coefficient polynomials times selected-variable dressings \(g\).
The already proved identities for \(g\) give
Equation~\eqref{eq:root-charge-words}.
\end{proof}

\subsection{Indecomposables and central moment values}

We record the geometric consequence of the charges in
Lemma~\ref{lem:root-charge-generation}.  For a dimension vector
\(\boldsymbol\gamma\), put
\[
 G_{\boldsymbol\gamma}=\prod_{i\in\Gamma_0}\GL_{\gamma_i},
 \qquad
 N_{\boldsymbol\gamma}=\bigoplus_{a:i\to j}
                     \Hom(\C^{\gamma_i},\C^{\gamma_j}).
\]
For a real vector \(\boldsymbol\theta=(\theta_i)\), let
\(\tau_{\boldsymbol\theta}\in\Lie(G_{\boldsymbol\gamma})^*\)
be the central functional
\[
 \tau_{\boldsymbol\theta}(\xi)
       =\sum_i\theta_i\operatorname{tr}(\xi_i).
\]
Choose standard Hermitian forms, and write
\((\mu_{\mathbb R},\mu_{\mathbb C})\) for the flat hyperkähler moment
maps on \(T^*N_{\boldsymbol\gamma}\).  We normalize the real parameter
so that an integral character \(\boldsymbol\theta\) corresponds to
real central level \(\tau_{\boldsymbol\theta}\) in the affine
Kempf--Ness correspondence.  With any other common normalization,
the same fixed scalar multiplies every level below.

\begin{lemma}\label{lem:root-moment}
If \(\boldsymbol\gamma\in\Delta_\Gamma\) and
\(\boldsymbol\theta\cdot\boldsymbol\gamma=0\), then
\[
 \mu_{\mathbb C}^{-1}(0)\cap
 \mu_{\mathbb R}^{-1}(\tau_{\boldsymbol\theta})\ne\varnothing.
\]
The scalar circle acting simultaneously on all nonzero vertex spaces
fixes every point of this representation.
\end{lemma}

\begin{proof}
First, \(\Gamma\) has an indecomposable complex representation \(x\)
of dimension \(\boldsymbol\gamma\).
For completeness, the required passage from Kac support to characteristic
zero can be made by constructibility.  A representation is decomposable
over an algebraic closure exactly when its endomorphism algebra contains
an idempotent other than \(0\) and \(1\).  The equations for such
idempotents, together with these two open exclusions, define a
constructible incidence condition over \(\Z\).
Chevalley's theorem therefore gives a constructible locus of
geometrically indecomposable representations.  A nonzero Kac polynomial
gives points of this locus in infinitely many residue characteristics.
Its generic fiber cannot be empty, since a constructible subset of
\(\Spec\Z\) that omits the generic point is finite.  Base change of
the nonempty characteristic-zero locus to \(\C\) gives \(x\).

The following trace-annihilator argument is the lifting criterion of
Crawley--Boevey \cite[Lemma~3.1 and Theorem~3.3]{CrawleyBoevey01}; we
recall it before passing to the real moment level.
Every endomorphism \(e\) of \(x\) has the form \(c\id+n\), where \(n\)
is nilpotent.  Indeed, its generalized eigenspaces are subrepresentations,
so indecomposability permits only one eigenvalue.
Consequently
\[
 \tau_{\boldsymbol\theta}(e)
 =c\,\boldsymbol\theta\cdot\boldsymbol\gamma=0.
\]
The cotangent moment map with \(x\) fixed is the transpose of the
infinitesimal orbit map:
\[
 N_{\boldsymbol\gamma}^*\longrightarrow
       \Lie(G_{\boldsymbol\gamma})^*,
 \qquad
 y\longmapsto\bigl(\xi\longmapsto y(\xi x)\bigr).
\]
Its image is the annihilator of
\(\End_\Gamma(x)=\ker(\xi\mapsto\xi x)\).  Hence the complex moment
fiber at \(\tau_{\boldsymbol\theta}\) is nonempty, as is the fiber at
any common scalar multiple of that value.

This fiber is a closed invariant affine subset, since the moment value
is central.  The closure of any orbit in it contains an orbit closed
in the ambient affine representation.  By the zero-character affine
Kempf--Ness theorem, that orbit meets \(\mu_{\mathbb R}^{-1}(0)\);
see \cite[Theorem~6.1 and Corollary~6.2]{King94}.
Thus there is a point with hyperkähler moment value
\((0,\tau_{\boldsymbol\theta},0)\), using the real and imaginary
parts of the complex moment map and adjusting the common scalar to
the chosen normalization.

The quaternionic rotation on the flat cotangent representation commutes
with the compact gauge group and rotates its three moment maps.
A single rotation sends this value to
\((\tau_{\boldsymbol\theta},0,0)\)
\cite[\S3(D)]{HKLR87}.
This proves the asserted nonemptiness.  The simultaneous scalar acts
trivially on all arrow spaces and their duals, proving the last statement.
\end{proof}

\subsection{The framing vertex and complementary monopoles}

\begin{proposition}\label{prop:cocharacter-fixed}
Suppose that every \(\nu\)-semistable point of \(\mu^{-1}(0)\) is stable
and has trivial \(G\)-stabilizer.  Then the natural quotient map
\[
 A_F/\langle (A_F)_m:m\ne0\rangle
       \longrightarrow B_F,
\]
where \(m\) denotes \(\nu\)-weight, is an isomorphism.
This is an equality of quotients of \(A_F\), preserving the coefficient
map and both gradings.
\end{proposition}

\begin{proof}
Lemma~\ref{lem:trivial-summands} allows us to suppose that \(F\) scales
the individual arrows and framing rows.  Adding or removing the
gauge-trivial cotangent summands in that lemma preserves both regularity
and all the fixed quotients under consideration.
Write
\[
 \nu((g_i))=\prod_i\det(g_i)^{\nu_i}.
\]
Adjoin one vertex \(\infty\) of dimension one and \(w_i\) arrows
\(i\to\infty\), keeping the original arrows.  Denote this unframed
quiver by \(\widehat Q\) and its dimension vector by
\(\mathbf d=(\mathbf v,1)\).
Its representation is the original \(N\), after a choice of basis
in each framing space.

There is a group isomorphism
\[
 G\times\C^\times\longrightarrow G_{\mathbf d},
 \qquad ((g_i),c)\longmapsto((cg_i),c).
\]
The second factor acts trivially on \(N\).  The BFN space of triples
and its convolution therefore split, giving
\begin{equation}\label{eq:central-splitting}
 A_{\widehat Q,\mathbf d,F}
       \cong A_F\otimes_\C\C[\xi,z,z^{-1}].
\end{equation}
Here \(\xi\) is the coefficient parameter of the extra scalar and
\(z\) is its fundamental Laurent monopole, of cohomological degree zero.
This \(z\) is unrelated to the degree-two added-loop mass in
Section~\ref{sec:deformations}.  In the original vertex
coordinates, \(z\) has topological charge \(\mathbf d\): it translates
every gauge coweight entry, including the entry at \(\infty\), by one.

Extend the character by
\[
 \widehat\nu_i=\nu_i\quad(i\in Q_0),
 \qquad
 \widehat\nu_\infty=-\sum_i\nu_i v_i.
\]
Then \(\widehat\nu\cdot\mathbf d=0\), and under the displayed group
isomorphism \(\widehat\nu\) is \(\nu\) on \(G\) and trivial on the
extra scalar.  In particular, both \(\xi\) and \(z^{\pm1}\) have
\(\widehat\nu\)-weight zero.  Taking the cocharacter fixed quotient
in Equation~\eqref{eq:central-splitting} consequently gives
\[
 \bigl(A_{\widehat Q,\mathbf d,F}\bigr)_{\mathrm{fix},\widehat\nu}
 \cong
 \bigl(A_F\bigr)_{\mathrm{fix},\nu}
                 \otimes_\C\C[\xi,z,z^{-1}],
\]
where the subscript denotes the quotient by nonzero cocharacter
weights.  This extension is faithfully flat.

Let \(M_{\mathbf b',p}^{+}\) be a nontrivial positive minuscule generator
of \(A_F\), and put \(\mathbf b=(\mathbf b',0)\).
If \(\widehat x_{ia}\) denotes an enlarged gauge Chern root, the group
isomorphism above identifies the original root with
\(x_{ia}=\widehat x_{ia}-\xi\).
Let \(p^{\uparrow}\) be the polynomial obtained from \(p\) by these
substitutions, leaving every flavor parameter unchanged.
The original generator is then the positive minuscule class
\(M_{\mathbf b,p^{\uparrow}}^{+}\).
Central translation also expresses it using the complementary negative
coweight:
\begin{equation}\label{eq:complementary-monopoles}
 M_{\mathbf b,p^{\uparrow}}^{+}
       =z\,M_{\mathbf d-\mathbf b,p^{\uparrow}}^{-}.
\end{equation}
The two coweights have the same Weyl stabilizer, so the same polynomial
\(p^{\uparrow}\) is an admissible dressing on both sides.
It need not depend only on the variables selected by the complementary
negative coweight; Lemma~\ref{lem:root-charge-generation} allows all
stabilizer-invariant dressings.
The equality follows directly from
Equation~\eqref{eq:minuscule}: all matter weights pair trivially
with the central coweight, and central translation identifies the
minuscule orbits and their equivariant tangent bundles.
Equivalently it follows from convolution with the invertible central
class.  Thus it holds with compatible fundamental-class conventions;
rescaling minuscule generators changes only a nonzero constant.

Suppose that this generator has nonzero image in the cocharacter
fixed quotient.  Apply Lemma~\ref{lem:root-charge-generation} in the
enlarged quiver to \(M_{\mathbf b,p^{\uparrow}}^{+}\).
Any word containing a factor of nonzero \(\widehat\nu\)-weight
vanishes in that quotient.  A nonzero image therefore requires a
decomposition
\[
 \mathbf b=\sum_a\boldsymbol\alpha_a,
 \qquad
 \boldsymbol\alpha_a\in\Delta_{\widehat Q},
 \qquad
 \widehat\nu\cdot\boldsymbol\alpha_a=0.
\]
Since \(z\) is an invertible element of weight zero,
Equation~\eqref{eq:complementary-monopoles} implies that
\(M_{\mathbf d-\mathbf b,p^{\uparrow}}^{-}\) also survives.  Applying the negative
half of the same lemma yields
\[
 \mathbf d-\mathbf b=\sum_b\boldsymbol\beta_b,
 \qquad
 \boldsymbol\beta_b\in\Delta_{\widehat Q},
 \qquad
 \widehat\nu\cdot\boldsymbol\beta_b=0.
\]
Both sums are nonempty: \(\mathbf b\ne0\), while its \(\infty\)-entry
is zero and that of \(\mathbf d\) is one.

Together these decompositions express \(\mathbf d\) as a sum of at
least two nonzero dimension vectors
\(\boldsymbol\gamma_1,\ldots,\boldsymbol\gamma_k\) in
\(\Delta_{\widehat Q}\), each orthogonal to \(\widehat\nu\).
For each summand choose the point furnished by
Lemma~\ref{lem:root-moment}, with real central parameter
\(\widehat\nu\).  Use orthogonal decompositions of the vertex spaces
and take the direct sum of these points.  Moment maps are block
diagonal on direct sums, so the resulting point has
\[
 \mu_{\mathbb C}=0,\qquad
 \mu_{\mathbb R}=\tau_{\widehat\nu}.
\]
Its stabilizer contains \(k\) independent block scalar circles.

The enlarged reduction is exactly the original framed reduction:
the extra diagonal scalar is ineffective, its moment equation is
redundant, and \(\widehat\nu\) is trivial on it.
After dividing by this scalar, the stabilizer still contains a torus
of dimension \(k-1>0\).
By the affine Kempf--Ness correspondence at the character level, the
resulting point is \(\nu\)-semistable in the original zero complex
moment fiber \cite[Theorem~6.1 and Corollary~6.2]{King94}.
This contradicts the assumed freeness of every such point.
The argument also applies when the stable locus is empty: it would
then produce a semistable point where none exists.

Every nontrivial positive minuscule generator therefore vanishes.
For a negative generator use the identity
\(M_{\mathbf b,p^{\uparrow}}^{-}
  =z^{-1}M_{\mathbf d-\mathbf b,p^{\uparrow}}^{+}\)
and interchange the two halves in the argument; it vanishes as well.
By Proposition~\ref{prop:minuscule}, the ideal they generate is the ideal
generated by the nonzero full topological-charge spaces.  It is therefore
contained in the ideal generated by the nonzero \(\nu\)-weight spaces.
For \(m\ne0\), the \(\nu\)-weight space \((A_F)_m\) is a sum of nonzero
topological-charge spaces.  This gives the opposite inclusion.
The ideals coincide, proving the proposition.
If all gauge ranks are zero there are no such generators and the
same conclusion is immediate.
\end{proof}

\subsection*{Completion of the proof}
\begin{proof}[Proof of Theorem~\ref{thm:main}]
Proposition~\ref{prop:cocharacter-fixed} and Theorem~\ref{thm:full-torus}
give graded $R$-algebra isomorphisms
\[
 \C[Y_{\mathfrak f}^{\nu}]
   \xrightarrow[\ \sim\ ]{\beta} B_F
   \xrightarrow[\ \sim\ ]{\alpha} H_F^*(X),
 \qquad (\alpha\circ\beta)\circ\phi_C=\phi_H.
\]
Thus $\ker\phi_C=\ker\phi_H$.  The inverse of $\alpha\circ\beta$ is the
isomorphism in Theorem~\ref{thm:main}; surjectivity of $\phi_H$ makes it the
unique isomorphism compatible with the coefficient maps.  All quotient
rings are scheme-theoretic and retain their nilpotents.
\end{proof}
\bibliographystyle{plain}
\bibliography{references}
\end{document}
